%% ============================================================================
%% Chapter 19: Testing & Quality Engineering for APIs
%% Mastering APIs: From Zero to Production Guru
%% ============================================================================
\chapter{Testing \& Quality Engineering for APIs}
\label{ch:testing}

%% ----------------------------------------------------------------------------
%% 1. Executive Summary
%% ----------------------------------------------------------------------------
An API you have not tested is a rumor. Worse, an API tested only by its own
developers is a rumor the provider tells itself: every chapter of this book
so far has built artifacts --- HTTP semantics (Chapter~\ref{ch:http_wire}),
serializers (Chapter~\ref{ch:serialization}), FastAPI services
(Chapter~\ref{ch:fastapi}), RFC~9457 error shapes
(Chapter~\ref{ch:problem_details}), versioned evolvable contracts
(Chapter~\ref{ch:versioning}), resilience policies
(Chapter~\ref{ch:microservices}) --- and each artifact makes promises that
\emph{someone else} will depend on. This chapter is about turning those
promises into executable evidence, and about building the quality gates that
refuse to ship when the evidence goes missing.

We begin with the \textbf{API test pyramid}: unit tests for validators and
serializers, component tests that run an endpoint against real in-process
dependencies, contract tests that pin the consumer--provider agreement,
integration tests that cross process boundaries, and a thin crust of
end-to-end journeys. For each layer we state precisely which defect classes
it can and cannot catch, its cost and speed profile, and where the
``testing trophy'' critique of the classic pyramid is right and wrong.
We then go deep on \textbf{consumer-driven contract testing} with the Pact
model~\cite{pactdocs}: the consumer records interactions into a pact file,
a broker distributes it, the provider replays it against real code using
named provider states, and a \texttt{can-i-deploy} gate blocks incompatible
releases. We contrast this with OpenAPI-diff schema checking
\cite{openapi31} and give a decision rule for when each applies.

Next come the two most under-used techniques in API engineering.
\textbf{Property-based testing} replaces hand-picked examples with
generated universes of inputs and machine-checked invariants --- never
5xx, schema-conformant 2xx, problem-shaped 4xx --- with automatic
\emph{shrinking} of counterexamples and rule-based \emph{stateful} testing
of API sequences. \textbf{Offline-first mocking} gives us a precise
fake/stub/mock taxonomy, \texttt{httpx.MockTransport}-style interception,
record--replay cassettes, and --- critically --- \emph{virtual clocks} that
let us test retries, timeouts, and circuit breakers deterministically
without ever calling \texttt{sleep}. Then we build \textbf{load and
performance testing} from first principles: open versus closed workload
models, coordinated omission, a derivation and application of Little's Law
$L = \lambda W$, saturation sweeps and knee-point analysis, and the
load/stress/soak/spike taxonomy \cite{nygard2018releaseit,beyer2016sre}.
\textbf{Chaos engineering} follows: a fault taxonomy (latency, errors,
aborts, partitions, resource exhaustion), blast-radius control, and
steady-state hypotheses, implemented as deterministic fault-injection
middleware. Finally we assemble everything into \textbf{CI quality gates as
code}: coverage floors, contract gates, performance regression budgets,
fuzz smoke tests, and disciplined flake management.

Every technique is backed by runnable, offline, deterministic Python~3.11
code against the synthetic Northwind Robotics robot registry: a mini
contract-testing framework that \emph{proves} it catches a breaking field
rename, a Hypothesis suite plus a stdlib OpenAPI fuzzer that finds and
shrinks a planted bug, an \texttt{asyncio} load harness that prints a
knee-point table, fault-injection middleware exercised entirely on virtual
time, and a single \texttt{ci\_gate.py} that runs the whole battery and
exits 0 or 1.

\begin{keyconceptbox}
A test suite is not a pile of checks; it is a \emph{supply chain of
evidence}. Each layer manufactures a different kind of confidence at a
different unit cost: units catch logic bugs in milliseconds, contracts
catch cross-team breakage without a shared environment, fuzzers explore
input spaces humans cannot enumerate, load harnesses locate the saturation
knee before your customers do, and chaos experiments prove your resilience
claims instead of assuming them. The craft of quality engineering is
allocating budget across these layers so that the \emph{cheapest} layer
that can catch each defect class is the one that does.
\end{keyconceptbox}

%% ----------------------------------------------------------------------------
\section{Learning Objectives \& Prerequisites}
\label{sec:ch19-objectives}

\subsection*{Learning Objectives}
After completing this chapter, its lab, and its exercises, you will be able to:

\begin{enumerate}[label=\textbf{LO-\arabic*.}, leftmargin=2.6cm]
  \item Place any given test on the API test pyramid (unit, component,
        contract, integration, end-to-end), state which defect classes that
        layer catches, and defend the cost/speed trade-off of the placement,
        including the testing-trophy critique of the classic pyramid.
  \item Implement the consumer-driven contract loop end to end: record
        consumer interactions into a pact file, distribute it, replay it on
        the provider against real code using named provider states, and
        gate deployment on verification --- and articulate when Pact-style
        contracts beat OpenAPI-diff schema checking and vice
        versa~\cite{pactdocs,openapi31}.
  \item Write property-based API tests with explicit strategies and
        invariants (never 5xx; 2xx validates against the response schema;
        4xx is RFC~9457 shaped), explain how shrinking minimizes
        counterexamples, and design rule-based stateful tests that catch
        sequence-dependent bugs.
  \item Distinguish fakes, stubs, mocks, and record--replay cassettes with
        precise definitions, and test retries, timeouts, and circuit
        breakers deterministically using virtual clocks and scripted fault
        injection instead of wall-clock sleeps~\cite{httpxdocs}.
  \item Build an offline load harness implementing open and closed workload
        models, derive and apply Little's Law $L = \lambda W$ to capacity
        questions, run a saturation sweep, and read a knee-point table to
        set performance budgets.
  \item Define a fault taxonomy (latency, error, abort, partition, resource
        exhaustion), state a steady-state hypothesis, bound a blast radius,
        and implement deterministic fault-injection middleware for tests.
  \item Assemble CI quality gates as code --- coverage floors, contract
        verification, fuzz smoke, performance smoke with structural (not
        wall-clock) assertions --- and operate a flake-management policy
        that keeps gates trustworthy.
\end{enumerate}

\subsection*{Prerequisites}
This chapter assumes you have completed:

\begin{itemize}
  \item \textbf{Chapters 0--5} --- Python~3.11+ environment, HTTP wire
        semantics, consuming APIs with \texttt{httpx}, serialization
        contracts, and client-side security hygiene (no credentials in
        tests, ever).
  \item \textbf{Chapter 6} (Building APIs with FastAPI) --- application
        factories, dependency injection, Pydantic models, and running an
        app in-process through an ASGI transport~\cite{fastapidocs}.
  \item \textbf{Chapter 7} (Errors, Validation \& Problem Details) ---
        RFC~9457 \texttt{application/problem+json} shape; our fuzz gate
        asserts it for every 4xx.
  \item \textbf{Chapter 9} (Versioning \& Evolution) --- additive versus
        breaking change; the contract-diff gate connects directly to the
        compatibility rules defined there.
  \item \textbf{Chapter 15} (Microservices \& Resilience) --- retries,
        timeouts, and circuit breakers, which we now learn to test
        deterministically.
  \item \textbf{Chapter 18} (Observability \& SRE) --- latency percentiles
        and SLIs; the load harness histogram reuses that chapter's
        percentile machinery.
\end{itemize}

All code runs offline on Windows with PowerShell: services are exercised
in-process via ASGI transports, pact files live on local disk, and every
random process is seeded or derandomized. Reported performance numbers are
illustrative, never asserted.

%% ----------------------------------------------------------------------------
\section{Deep Technical Mechanics \& Architectural Concepts}
\label{sec:ch19-mechanics}

\subsection{The API Test Pyramid}
\label{sec:ch19-pyramid}

The test pyramid is a portfolio theory, not a geometry lesson. Each layer
answers a different question, at a different speed, at a different cost per
defect found:

\begin{definition}[Test layers for APIs]
\label{def:test-layers}
A \textbf{unit test} exercises a pure unit of logic --- a validator, a
serializer, a pagination-token codec --- with no I/O. A \textbf{component
test} exercises one whole service in-process: real routing, real
validation, real serialization, with external dependencies replaced or
embedded. A \textbf{contract test} checks that a provider still satisfies
the recorded expectations of its consumers (or that a spec change is
backward compatible). An \textbf{integration test} crosses a real process
or infrastructure boundary: service to database, service to queue, service
to service. An \textbf{end-to-end (E2E) test} drives a user journey through
the deployed system, front door to database and back.
\end{definition}

\begin{figure}[htbp]
\centering
\begin{tikzpicture}[x=1cm, y=1cm,
  layer/.style={draw=packtgray, thick},
  lbl/.style={align=center, font=\small},
  side/.style={align=left, font=\footnotesize, text=packtgray, text width=6.6cm},
  axis/.style={->, thick, packtgray}]
  % --- pyramid body (trapezoids, bottom widest) ---
  \draw[layer, fill=blue!7]   (-5.2,0.0) -- (5.2,0.0) -- (4.0,1.4) -- (-4.0,1.4) -- cycle;
  \draw[layer, fill=green!9]  (-4.0,1.4) -- (4.0,1.4) -- (2.8,2.8) -- (-2.8,2.8) -- cycle;
  \draw[layer, fill=orange!12] (-2.8,2.8) -- (2.8,2.8) -- (1.6,4.2) -- (-1.6,4.2) -- cycle;
  \draw[layer, fill=purple!10] (-1.6,4.2) -- (1.6,4.2) -- (0.6,5.4) -- (-0.6,5.4) -- cycle;
  \draw[layer, fill=red!10]   (-0.6,5.4) -- (0.6,5.4) -- (0.0,6.3) -- cycle;
  \node[lbl] at (0,0.68) {\textbf{Unit} --- validators, serializers, codecs, pure rules};
  \node[lbl] at (0,2.08) {\textbf{Component} --- endpoint + real deps, in-process ASGI};
  \node[lbl] at (0,3.48) {\textbf{Contract} --- pact verify, schema diff};
  \node[lbl] at (0,4.78) {\textbf{Integration}};
  \node[lbl] at (0,5.72) {\textbf{E2E}};
  % --- side annotations ---
  \node[side, anchor=west] at (6.0,0.68) {thousands of cases; microseconds each; pinpoints the broken line};
  \node[side, anchor=west] at (6.0,2.08) {hundreds of cases; milliseconds each; catches wiring + validation + error-shape bugs};
  \node[side, anchor=west] at (6.0,3.48) {one suite per consumer--provider pair; catches cross-team breakage without a shared environment};
  \node[side, anchor=west] at (6.0,4.78) {tens of cases; seconds each; catches config, auth, schema-migration drift};
  \node[side, anchor=west] at (6.0,5.72) {a handful of golden journeys; minutes each; proves the system, not the parts};
  % --- axes ---
  \draw[axis] (-6.3,0.1) -- (-6.3,6.1) node[above, note, align=center] {cost, time,\\brittleness};
  \draw[axis] (-6.9,6.1) -- (-6.9,0.1) node[below, note, align=center] {execution\\speed};
\end{tikzpicture}
\caption{The API test pyramid. Width encodes recommended \emph{volume};
height encodes cost and environmental fidelity. Contract tests sit above
component tests because they assume a working provider component, yet far
below integration tests in environmental cost.}
\label{fig:ch19-pyramid}
\end{figure}

The economics are unforgiving. A defect caught by a unit test costs a
millisecond rerun; the same defect caught by an E2E test costs a
distributed-systems debugging session across three teams' dashboards. The
inverted pyramid (``ice-cream cone'') --- a thin unit layer and a fat E2E
layer --- is the single most common structural failure of API test suites:
it is slow, flaky, and worst of all \emph{ambiguous}, because a red E2E
test tells you that \emph{something} is wrong somewhere in five services.

\begin{remark}[The testing trophy debate]
Kent C.~Dodds' ``testing trophy'' argues the pyramid under-weights
\emph{integration} (in our terms, component) tests: modern in-process
transports and embedded fakes make whole-service tests cheap enough to
carry most of the suite. Both shapes agree on the actual invariant: push
verification \emph{down} to the cheapest layer that can catch the defect,
and keep the top thin. For HTTP APIs specifically, the component layer is
extraordinarily powerful because ASGI transports let us exercise the entire
service --- routing, validation, serialization, error shaping --- with zero
network, which is why most listings in this chapter live there.
\end{remark}

A useful discipline is to declare, per defect class, the \emph{designated
catcher}: validation bugs belong to unit and property tests; wiring and
serialization bugs to component tests; cross-team contract drift to
contract tests; infrastructure misconfiguration to integration tests; and
only genuine emergent, whole-system behavior to E2E. When a bug escapes to
production, the postmortem question is never ``why did no test catch it''
but ``which designated catcher failed, and why''.

\subsection{Consumer-Driven Contract Testing}
\label{sec:ch19-contracts}

Schema tells you what a provider \emph{can} serve; a consumer-driven
contract tells you what consumers \emph{actually rely on}. The gap between
those two sets is where deploys go to die: a provider renames a field it
believes unused, every one of its own tests stays green, and three consumer
services fail in production.

\begin{definition}[Consumer-driven contract]
\label{def:cdc}
A \textbf{consumer-driven contract} is a machine-readable artifact,
authored from the consumer's perspective, recording the interactions the
consumer depends on --- request shape, expected status, and the subset of
the response the consumer reads --- together with the \emph{provider
states} under which those expectations hold. The provider's build
\emph{verifies} the contract by replaying every interaction against real
provider code.
\end{definition}

\begin{definition}[Provider state]
\label{def:provider-state}
A \textbf{provider state} is a named precondition (``a robot with id 7
exists'') that the provider must arrange before replaying an interaction.
Good states are \emph{minimal} (they arrange exactly what the interaction
needs), \emph{independent} (any order, any subset), and \emph{semantic}
(named by domain meaning, not by fixture mechanics). A state that seeds the
entire universe is a smell: it hides which interactions need which data and
makes failures undebuggable.
\end{definition}

\begin{figure}[htbp]
\centering
\begin{tikzpicture}[x=1cm, y=1cm, node distance=0.9cm,
  box/.style={flowstep, fill=blue!5},
  artifact/.style={flowstep, fill=orange!15},
  gate/.style={flowstep, fill=red!8},
  wire/.style={->, thick},
  back/.style={->, thick, dashed, packtgray}]
  \node[box] (consumer) {\textbf{Consumer test}\\records interactions};
  \node[artifact, right=of consumer] (pact) {\textbf{pact file}\\consumer + provider\\+ interactions};
  \node[artifact, right=of pact] (broker) {\textbf{Broker}\\versioned contract\\store};
  \node[box, right=of broker] (provider) {\textbf{Provider CI}\\replay + provider\\states};
  \node[gate, below=0.9cm of provider] (candeploy) {\textbf{can-i-deploy}\\all pacts verified\\for this version?};
  \node[box, left=of candeploy] (deploy) {\textbf{Deploy / block}};
  \draw[wire] (consumer) -- node[above, note]{1. publish} (pact);
  \draw[wire] (pact) -- node[above, note]{2. store} (broker);
  \draw[wire] (broker) -- node[above, note]{3. fetch} (provider);
  \draw[wire] (provider) -- node[right, note]{4. results} (candeploy);
  \draw[wire] (candeploy) -- node[above, note]{5. gate} (deploy);
  \draw[back] (provider.north) .. controls +(0,1.1) and +(0,1.1) .. node[above, note]{6. verification results flow back to consumers} (consumer.north);
\end{tikzpicture}
\caption{The consumer-driven contract loop (Pact model~\cite{pactdocs}).
The pact file is written by the \emph{consumer}, owned by \emph{both}
teams, and executed by the \emph{provider's} build.}
\label{fig:ch19-contract-loop}
\end{figure}

Three properties make the loop work in organizations, not just in
diagrams. First, \textbf{the artifact crosses team boundaries}: the pact
file is a message from the consumer team to the provider team, and the
broker makes that message durable, versioned, and queryable. Second,
\textbf{verification runs on the provider's real code}, typically
in-process against the real application with real serialization --- a
contract verified against a mock of the provider verifies nothing. Third,
\textbf{deployment is gated}: \texttt{can-i-deploy} answers ``does version
$v$ of the provider satisfy every pact published against it?'' before $v$
ships, per environment.

\paragraph{Pact versus OpenAPI-diff.}
An OpenAPI-diff gate (connecting to Chapter~\ref{ch:versioning}) compares
two spec revisions and flags structural incompatibilities: removed
operations, removed response fields, tightened constraints. It needs no
consumer participation and works for public APIs with thousands of unknown
consumers --- but it can only flag what is \emph{structurally} breaking,
and it must conservatively treat every published field as load-bearing.
Pact-style contracts are usage-based: a field nobody reads can be renamed
without tripping any pact, and a \emph{semantically} breaking change inside
a structurally identical shape (same field, new meaning) can still be
caught by an explicit consumer expectation. The decision rule:

\begin{center}
\begin{tabularx}{\textwidth}{l X X}
\toprule
 & \textbf{Contract tests (Pact)} & \textbf{OpenAPI diff} \\
\midrule
Consumers & few, known, internal & many / unknown / public \\
Detects & usage-level breakage, semantic drift & structural breakage only \\
Setup cost & broker, provider states, both teams & spec in repo, diff tool \\
False safety & uncovered consumers & ``compatible'' spec, changed meaning \\
\bottomrule
\end{tabularx}
\end{center}

Mature platforms run both: schema diff as the coarse public-safety net,
consumer contracts as the precise internal gate.

\subsection{Property-Based Testing and API Fuzzing}
\label{sec:ch19-pbt}

Example-based tests enumerate inputs the author imagined; property-based
tests enumerate inputs the author \emph{did not}. A property-based test
consists of three parts: a \textbf{strategy} that generates a universe of
inputs (valid \emph{and} invalid), an \textbf{invariant} that must hold for
every generated input, and a \textbf{shrinker} that, on failure,
automatically minimizes the counterexample to the smallest reproducing
case. For HTTP APIs, the killer invariants are remarkably uniform:

\begin{enumerate}
  \item \textbf{Never 5xx.} Whatever the payload --- wrong types, missing
        fields, boundary values, pathological lengths --- the server must
        answer 2xx or a well-formed 4xx. A 5xx is an unhandled exception,
        full stop.
  \item \textbf{2xx validates against the published schema.} If the
        response does not conform to the contract you advertised, your
        successful responses are bugs.
  \item \textbf{4xx is RFC~9457 shaped.} Every client-side error carries
        \texttt{type}, \texttt{title}, \texttt{status}, \texttt{detail},
        \texttt{instance} (Chapter~\ref{ch:problem_details}), so consumers
        can handle errors programmatically.
\end{enumerate}

\textbf{Shrinking} is what makes failures actionable. A naive fuzzer
reports ``500 on payload of 4{,}000 random keys''; a shrinker
systematically deletes keys and minimizes values while the failure
reproduces, converging to something like
\texttt{\{"price\_cents": 666\}} --- which is a bug report a human can act
on. Hypothesis performs this automatically; our mini fuzzer implements a
greedy version by hand so the mechanics are visible.

\textbf{Stateful (rule-based) testing} extends properties across
\emph{sequences}: generate random interleavings of create/read/update
operations, maintain a local model of expected state, and assert the API
agrees with the model after every step. This catches the bugs single-call
tests structurally cannot: lost updates, read-your-writes violations,
idempotency-key reuse bugs (Chapter~\ref{ch:idempotency}), and leakage
between resources.

\textbf{Schemathesis-style fuzzing} derives the strategy automatically
from the OpenAPI document: every operation, every parameter, boundary and
malformed values per schema constraint, executed against the live service
with the same three invariants above~\cite{openapi31}. In
Section~\ref{sec:ch19-code} we build a stdlib miniature of this pipeline so
that the generation logic --- candidate values at \texttt{minLength}$-1$,
\texttt{maxLength}$+1$, \texttt{minimum}$-1$, wrong types, \texttt{null},
missing required keys, extra keys, non-object bodies --- is explicit rather
than magical.

\subsection{Test Doubles, Record--Replay, and Virtual Clocks}
\label{sec:ch19-doubles}

The word ``mock'' is used for five different things; imprecision here
causes real design damage, so we fix the taxonomy (after Meszaros'
xUnit patterns):

\begin{definition}[Test doubles]
\label{def:doubles}
A \textbf{dummy} is passed to satisfy a signature and never used. A
\textbf{stub} returns canned answers to calls the test triggers. A
\textbf{fake} is a working implementation with shortcuts (an in-memory
repository, an ASGI in-process transport). A \textbf{spy} records calls so
the test can assert on them afterward. A \textbf{mock} is programmed with
\emph{expectations} up front and fails on unexpected calls.
\end{definition}

For HTTP clients, the practical hierarchy is: prefer the \emph{fake}
(in-process ASGI transport to the real app) for anything you own; use
\emph{stubs} (\texttt{httpx.MockTransport} handlers, or \texttt{respx}
routes) for third-party edges with fixed canned responses; use
\textbf{record--replay cassettes} (VCR-style: record real interactions
once, replay byte-identical responses forever after) when the third party
is too complex to hand-stub --- with two disciplines: cassettes must be
\emph{sanitized} (no secrets, no real PII --- synthetic Northwind Robotics
data only) and \emph{re-recorded on a schedule}, because a cassette is a
contract claim that silently rots. WireMock-style standalone HTTP servers
are the same stub idea hoisted into a separate process for polyglot test
suites. The classic anti-pattern is mocking \emph{what you do not own} and
then asserting on the mock's call list: you have then tested that your code
calls your imagination of the vendor's API, and the test stays green while
the real integration rots.

\textbf{Virtual clocks} solve the hardest determinism problem in client
testing: time. Any test that calls \texttt{time.sleep(2)} to exercise a
retry backoff is slow, flaky, and unscalable. The fix is to inject the
sleeper: production wires \texttt{asyncio.sleep}; tests wire a
\texttt{VirtualClock} that records the requested delay and advances a
counter instantly. The test then asserts on the \emph{recorded schedule}
(\texttt{[0.1, 0.2, 0.4]} for exponential backoff) --- which is the actual
behavioral contract --- in microseconds of real time. The same injection
point serves timeouts and circuit breakers: fault injection decides
\emph{what} fails, the virtual clock decides \emph{how long} anything
takes, and no assertion ever touches the wall clock.

\subsection{Load and Performance Testing}
\label{sec:ch19-load}

Performance testing has one question: \emph{where is the knee, and how much
headroom separates us from it?} Everything else --- tools, models, charts
--- serves that question.

\paragraph{Open versus closed workload models.}
In a \textbf{closed model}, a fixed population of $N$ workers each issues a
request, waits for the response, then issues the next; offered load adapts
to the system's speed, so an overloaded system sees \emph{less} load ---
which hides overload. In an \textbf{open model}, requests arrive at a fixed
rate $\lambda$ regardless of completions; overload manifests as queueing
and latency growth, which is what real traffic does to real APIs. k6
(via arrival-rate executors) and Locust (via fixed user counts versus
custom arrival control) embody the same distinction. The classic
methodological sin is \textbf{coordinated omission}: a closed-model
load generator that schedules by \emph{completion} time silently skips the
load that should have been offered during a latency spike, and reports a
healthier system than reality.

\paragraph{Little's Law.}
\begin{theorem}[Little's Law]
\label{thm:little}
For a stable system, the mean number of requests in residence $L$, the
arrival rate $\lambda$, and the mean residence time $W$ satisfy
$L = \lambda W$.
\end{theorem}

\begin{proof}[Derivation]
Observe the system over a long interval $[0, T]$ in which $N$ requests
complete and the number in residence $L(t)$ returns to its initial value.
Each request $i$ contributes exactly its residence time $W_i$ to the area
under the curve $L(t)$:
\begin{align}
\int_{0}^{T} L(t)\,dt &= \sum_{i=1}^{N} W_i.
\end{align}
Divide both sides by $T$:
\begin{align}
L \;=\; \frac{1}{T}\int_{0}^{T} L(t)\,dt \;=\; \frac{N}{T}\cdot\frac{1}{N}\sum_{i=1}^{N} W_i \;=\; \lambda \cdot W,
\end{align}
since $\lambda = N/T$ is the throughput (equal to the arrival rate in a
stable system) and $W = \frac{1}{N}\sum_i W_i$ is the mean residence time.
\end{proof}

Little's Law is the Swiss army knife of capacity reasoning. If your API
serves $\lambda = 500$ req/s at a mean latency of $W = 40$ ms, then on
average $L = 500 \times 0.04 = 20$ requests are in flight --- your worker
pool, connection pool, and async concurrency limits must all exceed 20 with
headroom, or queueing begins. Conversely, in a closed-model test with $N$
workers you can \emph{predict} throughput: $\lambda = N / W$. Our harness
performs exactly this cross-check numerically.

\paragraph{Saturation and the knee.}
Let $S$ be mean service time and $c$ the effective concurrency (workers,
event-loop capacity, pool size). Utilization is $\rho = \lambda S / c$, and
queueing theory gives the qualitative shape $W \approx S / (1 - \rho)$:
latency is flat while $\rho \ll 1$, then explodes hyperbolically as
$\rho \to 1$. The \textbf{knee} is the offered rate at which the curve
turns --- operationally, the first load level where tail latency degrades
superlinearly (we flag p99 $>$ 10$\times$ baseline) or errors appear.
Sustainable capacity is provisioned \emph{below} the knee, typically at
60--70\% of it, because arrival bursts and deploys eat the remainder
\cite{beyer2016sre,nygard2018releaseit}.

\paragraph{The four load test shapes.}
\begin{center}
\begin{tabularx}{\textwidth}{l X l}
\toprule
\textbf{Shape} & \textbf{Question answered} & \textbf{Duration} \\
\midrule
Load   & Does the system meet SLOs at \emph{expected} peak? & 15--60 min \\
Stress & Where is the knee? What fails first past it? & ramp to failure \\
Soak   & Do leaks, fragmentation, or log growth degrade us over time? & hours--days \\
Spike  & Does the system survive a 10$\times$ burst (cache cold, autoscaler lag)? & seconds--minutes \\
\bottomrule
\end{tabularx}
\end{center}

Soak tests catch what ramps cannot: memory leaks, file-descriptor
exhaustion, token-cache stampede alignment, and disk growth
(Chapter~\ref{ch:observability} covers the telemetry). Spike tests expose
\emph{elasticity} defects --- the load balancer scales in minutes, the
traffic arrives in seconds.

\begin{figure}[htbp]
\centering
\begin{tikzpicture}[x=1cm, y=1cm, node distance=0.55cm,
  box/.style={flowstep, fill=green!7},
  analysis/.style={flowstep, fill=orange!12},
  gate/.style={flowstep, fill=red!8},
  wire/.style={->, thick}]
  \node[box] (model) {\textbf{1. Model workload}\\open or closed?\\request mix, payload sizes};
  \node[box, right=of model] (baseline) {\textbf{2. Baseline}\\low-rate run\\establishes p50 reference};
  \node[box, right=of baseline] (sweep) {\textbf{3. Saturation sweep}\\ramp offered RPS\\fixed short windows};
  \node[analysis, right=of sweep] (knee) {\textbf{4. Knee analysis}\\p99 vs baseline,\\error rate, Little check};
  \node[analysis, below=0.55cm of knee] (budget) {\textbf{5. Budgets}\\capacity headroom\\perf regression limits};
  \node[gate, left=of budget] (ci) {\textbf{6. CI perf gate}\\structural assertions\\smoke sweep offline};
  \draw[wire] (model) -- (baseline);
  \draw[wire] (baseline) -- (sweep);
  \draw[wire] (sweep) -- (knee);
  \draw[wire] (knee) -- (budget);
  \draw[wire] (budget) -- (ci);
  \draw[wire, dashed, packtgray] (ci.west) -- ++(-1.1,0) |- node[left, note, align=center]{re-calibrate on\\arch change} (model.west);
\end{tikzpicture}
\caption{Load-test methodology. The CI gate runs a short structural smoke;
deep sweeps are scheduled offline and re-baselined on architectural
change.}
\label{fig:ch19-load-flow}
\end{figure}

\subsection{Chaos Engineering Basics}
\label{sec:ch19-chaos}

Chaos engineering is the experimental validation of resilience claims.
Chapter~\ref{ch:microservices} taught you to \emph{build} retries,
timeouts, and breakers; chaos asks: do they actually work, in combination,
under realistic failure?

\begin{definition}[Steady-state hypothesis and blast radius]
\label{def:chaos}
The \textbf{steady state} is a measurable, behavioral definition of
``the system is fine'' --- e.g., p99 $<$ 250 ms and error rate $<$ 0.1\%
over five minutes --- never an internal one (``CPU looks low''). A
\textbf{steady-state hypothesis} asserts the system remains in steady state
under a specified fault. The \textbf{blast radius} is the maximal scope the
experiment can affect; experiments start at the smallest blast radius (one
in-process middleware in a test) and widen only as confidence grows.
\end{definition}

\begin{center}
\begin{tabularx}{\textwidth}{l X X}
\toprule
\textbf{Fault class} & \textbf{Mechanism} & \textbf{Reveals} \\
\midrule
Latency & inject delay before/around the dependency & missing timeouts, thread/connection starvation \\
Error (5xx) & synthesize error responses & retry storms, missing backoff/jitter, error-mapping bugs \\
Abort & drop/reset the connection with no response & retry-on-exception paths, hung-read behavior \\
Partition & make a dependency unreachable entirely & failover, cache fallback, graceful degradation \\
Resource exhaustion & cap CPU/memory/threads/FDs & queue growth, priority inversion, GC death spirals \\
\bottomrule
\end{tabularx}
\end{center}

Distributed systems fail in partial, correlated, and Byzantine ways
\cite{kleppmann2017ddia}; the taxonomy above is ordered roughly by
increasing subtlety. The critical engineering move for \emph{testing} is
determinism: our fault-injection middleware is scripted by request index,
not by random dice, so a failing resilience test reproduces byte-for-byte.
Randomized chaos is for staging environments and game days; scripted chaos
is for CI.

\subsection{CI Quality Gates as Code}
\label{sec:ch19-gates}

A quality gate is a function from change to evidence, executed on every
pull request, whose failure blocks merge. The design rules that keep gates
healthy:

\begin{enumerate}
  \item \textbf{Hermetic.} No network, no cloud, no shared environment.
        Everything in this chapter runs in-process; a gate that needs
        staging is a gate that will be disabled.
  \item \textbf{Fast.} Minutes, not hours. Deep load sweeps and long soaks
        are scheduled jobs that update budgets; the PR gate runs structural
        smoke.
  \item \textbf{Self-checking.} A gate that cannot fail proves nothing.
        Our gate runs each check against a deliberately broken variant and
        requires the check to catch it --- mutation testing applied to the
        pipeline itself.
  \item \textbf{Deterministic.} Derandomized property seeds, scripted
        faults, structural perf assertions. A flaky gate trains the team to
        click ``re-run'' without reading, which is worse than no gate.
  \item \textbf{Legible.} The report names the layer, the interaction, and
        the mismatch: \texttt{\$.sku: missing}. ``1 test failed'' is not a
        diagnosis.
\end{enumerate}

Typical API gate battery: unit + property suites with a coverage floor
(floors guard against suite deletion, not for their own sake --- an
80\% floor with the invariants of Section~\ref{sec:ch19-pbt} beats a 95\%
floor of example tests); contract verification against every published
pact plus an OpenAPI diff against the previous released spec; a fuzz smoke
(a bounded case budget, not an open-ended campaign); a perf smoke that
asserts completion and zero 5xx rather than latency (CI machines are noisy
neighbors; latency budgets belong to scheduled runs on quiet hardware);
and flake management: quarantine list with SLA, automatic bisection of
derandomized failures, and a hard rule that a test flaky twice in a week is
either fixed or deleted. Delivery-performance research consistently finds
that fast, trusted automated gates correlate with both throughput and
stability --- the gate is a delivery feature, not a tax
\cite{beyer2016sre}.

%% ----------------------------------------------------------------------------
\section{Production-Ready Code Implementation}
\label{sec:ch19-code}

All listings live under \texttt{code\textbackslash{}ch19\textbackslash{}}
and share one app under test: the Northwind Robotics robot registry. They
are fully offline (in-process ASGI), deterministic (derandomized
generators, scripted faults, virtual clocks), and use synthetic data only.
Install dependencies once:

\begin{minted}{powershell}
cd code\ch19
pip install -r requirements.txt
\end{minted}

\subsection{Listing 19.1 --- The App Under Test with Deliberate-Defect Knobs}
\label{sec:ch19-listing-app}

The registry (Listing 19.1) is a plain Chapter~6 FastAPI service
with two teaching additions. First, every error is RFC~9457 problem-shaped
(Chapter~\ref{ch:problem_details}), because later listings \emph{assert} on
that shape. Second, \texttt{create\_app} accepts two keyword-only defect
knobs: \texttt{breaking=True} renames \texttt{sku} to
\texttt{part\_number} in responses (a contract break), and
\texttt{bug\_at\_price=N} returns an unstructured 500 when
\texttt{price\_cents == N} (a planted bug for the fuzzer). The module also
provides \texttt{SyncASGITransport}: since \texttt{httpx} 0.28 the bundled
\texttt{ASGITransport} is async-only~\cite{httpxdocs}, so we wrap it in a
five-line shim that runs the app on one dedicated event loop and re-wraps
the body as a sync byte stream.

\begin{minted}{python}
"""Northwind Robotics robot registry: the single API under test in Chapter 19.

``create_app`` ships with two deliberate-defect knobs so the lab can *break*
the API on demand:

* ``breaking=True`` renames the ``sku`` field to ``part_number`` in every
  response body -- the classic contract-breaking change a provider team
  ships believing it is "internal".
* ``bug_at_price=N`` makes ``POST /robots`` return an unhandled 500 (not
  RFC 9457 shaped) whenever ``price_cents == N`` -- a planted bug for the
  fuzzer to find and shrink.

Offline-first: every consumer of this module talks to the app in-process
through ``httpx.ASGITransport``; no socket is ever opened.
"""
from __future__ import annotations

import asyncio
from typing import Any

import httpx
from fastapi import FastAPI, HTTPException, Request
from fastapi.exceptions import RequestValidationError
from fastapi.responses import JSONResponse
from pydantic import BaseModel, Field

PROBLEM_BASE = "https://northwind-robotics.example/problems"


class SyncASGITransport(httpx.BaseTransport):
    """Sync shim over ``httpx.ASGITransport`` (async-only since httpx 0.28).

    Runs the ASGI app on one dedicated event loop, in-process: no sockets,
    no network, fully offline.
    """

    def __init__(self, app: Any) -> None:
        self._inner = httpx.ASGITransport(app=app)
        self._loop = asyncio.new_event_loop()

    def handle_request(self, request: httpx.Request) -> httpx.Response:
        async_resp = self._loop.run_until_complete(
            self._inner.handle_async_request(request)
        )
        content = self._loop.run_until_complete(async_resp.aread())
        # Rebuild with a byte body so the sync client gets a SyncByteStream.
        return httpx.Response(
            status_code=async_resp.status_code,
            headers=async_resp.headers,
            content=content,
        )


def make_sync_client(app: Any, base_url: str = "http://testserver") -> httpx.Client:
    """Build a synchronous httpx client wired straight into the ASGI app."""
    return httpx.Client(transport=SyncASGITransport(app), base_url=base_url)


class RobotCreate(BaseModel):
    """Payload accepted by ``POST /robots``."""

    sku: str = Field(min_length=3, max_length=32)
    name: str = Field(min_length=1, max_length=80)
    price_cents: int = Field(ge=0, le=10_000_000)


class Robot(RobotCreate):
    """Stored representation: everything in ``RobotCreate`` plus a server id."""

    id: int


def _problem(status: int, title: str, detail: str, instance: str) -> JSONResponse:
    """Build an RFC 9457 problem-details response."""
    slug = title.lower().replace(" ", "-")
    return JSONResponse(
        status_code=status,
        media_type="application/problem+json",
        content={
            "type": f"{PROBLEM_BASE}/{slug}",
            "title": title,
            "status": status,
            "detail": detail,
            "instance": instance,
        },
    )


def create_app(*, breaking: bool = False, bug_at_price: int | None = None) -> FastAPI:
    """Application factory; defect knobs are keyword-only and default to off."""
    app = FastAPI(title="Northwind Robotics -- Robot Registry")
    app.state.db = {}  # type: dict[int, dict[str, Any]]
    app.state.next_id = 1

    @app.exception_handler(RequestValidationError)
    async def on_validation_error(
        request: Request, exc: RequestValidationError
    ) -> JSONResponse:
        errors = exc.errors()
        first = errors[0] if errors else {}
        loc = ".".join(str(part) for part in first.get("loc", []))
        detail = f"{loc}: {first.get('msg', 'invalid request')}"
        return _problem(422, "Validation Error", detail, request.url.path)

    @app.exception_handler(HTTPException)
    async def on_http_error(request: Request, exc: HTTPException) -> JSONResponse:
        title = "Not Found" if exc.status_code == 404 else "HTTP Error"
        return _problem(exc.status_code, title, str(exc.detail), request.url.path)

    def present(record: dict[str, Any]) -> dict[str, Any]:
        """Serialize a stored record, applying the deliberate break if asked."""
        body = dict(record)
        if breaking:
            body["part_number"] = body.pop("sku")
        return body

    @app.post("/robots", status_code=201)
    async def create_robot(payload: RobotCreate) -> Any:
        if bug_at_price is not None and payload.price_cents == bug_at_price:
            # Planted bug: an unstructured, non-RFC-9457 server error.
            return JSONResponse(status_code=500, content={"oops": "price demon"})
        robot_id = app.state.next_id
        app.state.next_id += 1
        record = {"id": robot_id, **payload.model_dump()}
        app.state.db[robot_id] = record
        return present(record)

    @app.get("/robots/{robot_id}")
    async def get_robot(robot_id: int) -> dict[str, Any]:
        if robot_id not in app.state.db:
            raise HTTPException(status_code=404, detail=f"robot {robot_id} not found")
        return present(app.state.db[robot_id])

    return app
\end{minted}

\subsection{Listing 19.2 --- A Mini Pact: Recorder, Pact File, Provider Verifier}
\label{sec:ch19-listing-pact}

Listing 19.2 implements the full consumer-driven loop of
Figure~\ref{fig:ch19-contract-loop} in about 200 lines. The consumer side
(\texttt{PactRecorder}) executes each interaction, fails the consumer test
immediately if the API already violates the expectation, and freezes the
interaction --- method, path, request body, expected status, and a
\emph{body subset} (what the consumer actually reads) --- into JSON. The
provider side (\texttt{ProviderVerifier}) builds a \emph{fresh} application
per interaction, applies the named provider state, replays the request
against real code in-process, and reports recursive subset mismatches with
JSONPath-style locations. \texttt{PROVIDER\_STATES} is the provider team's
half of the bargain: the implementation of the consumer's named
preconditions.

\begin{minted}{python}
"""A mini consumer-driven contract-testing framework, Pact-inspired.

The moving parts mirror the real Pact loop:

1. Consumer side: :class:`PactRecorder` wraps an ``httpx.Client``; every
   interaction the consumer relies on is executed, validated on the spot,
   and frozen into a JSON *pact file*.
2. The pact file is the contract artifact. In production it travels through
   a broker; here it is a plain JSON file under ``pacts\\``.
3. Provider side: :class:`ProviderVerifier` replays each interaction against
   a *fresh* provider application, applying the named *provider state*
   first, and checks status + a recursive body subset match.

Everything is offline: provider replay runs through ``httpx.ASGITransport``.
"""
from __future__ import annotations

import json
from dataclasses import dataclass, field
from pathlib import Path
from typing import Any, Callable

import httpx

from robot_registry import create_app, make_sync_client

StateHandler = Callable[[Any], None]  # receives the provider FastAPI app
AppFactory = Callable[[], Any]


@dataclass
class Interaction:
    """One request/response pair the consumer depends on."""

    description: str
    provider_state: str | None
    method: str
    path: str
    request_body: dict[str, Any] | None
    expected_status: int
    expected_body_subset: dict[str, Any]

    def to_json(self) -> dict[str, Any]:
        return {
            "description": self.description,
            "providerState": self.provider_state,
            "request": {"method": self.method, "path": self.path, "body": self.request_body},
            "response": {
                "status": self.expected_status,
                "bodySubset": self.expected_body_subset,
            },
        }

    @classmethod
    def from_json(cls, raw: dict[str, Any]) -> "Interaction":
        return cls(
            description=raw["description"],
            provider_state=raw.get("providerState"),
            method=raw["request"]["method"],
            path=raw["request"]["path"],
            request_body=raw["request"].get("body"),
            expected_status=raw["response"]["status"],
            expected_body_subset=raw["response"]["bodySubset"],
        )


def subset_mismatches(expected: Any, actual: Any, path: str = "$") -> list[str]:
    """Recursive subset match: every expected leaf must exist and be equal."""
    if isinstance(expected, dict):
        if not isinstance(actual, dict):
            return [f"{path}: expected object, got {type(actual).__name__}"]
        problems: list[str] = []
        for key, sub in expected.items():
            if key not in actual:
                problems.append(f"{path}.{key}: missing (expected {sub!r})")
            else:
                problems.extend(subset_mismatches(sub, actual[key], f"{path}.{key}"))
        return problems
    if isinstance(expected, list):
        if not isinstance(actual, list) or len(actual) != len(expected):
            return [f"{path}: expected list of length {len(expected)}"]
        problems = []
        for i, sub in enumerate(expected):
            problems.extend(subset_mismatches(sub, actual[i], f"{path}[{i}]"))
        return problems
    if expected != actual:
        return [f"{path}: expected {expected!r}, got {actual!r}"]
    return []


class PactRecorder:
    """Consumer side: exercise the API and freeze the interactions relied on."""

    def __init__(self, client: httpx.Client, *, consumer: str, provider: str) -> None:
        self._client = client
        self.consumer = consumer
        self.provider = provider
        self.interactions: list[Interaction] = []

    def record(
        self,
        description: str,
        method: str,
        path: str,
        *,
        provider_state: str | None = None,
        json_body: dict[str, Any] | None = None,
        expected_status: int,
        expect_subset: dict[str, Any],
    ) -> httpx.Response:
        """Execute one interaction; fail the consumer test immediately if the
        API already violates the expectation, then freeze it."""
        resp = self._client.request(method, path, json=json_body)
        if resp.status_code != expected_status:
            raise AssertionError(
                f"[{description}] expected {expected_status}, got "
                f"{resp.status_code}: {resp.text}"
            )
        problems = subset_mismatches(expect_subset, resp.json())
        if problems:
            raise AssertionError(f"[{description}] response mismatches: {problems}")
        self.interactions.append(
            Interaction(
                description=description,
                provider_state=provider_state,
                method=method,
                path=path,
                request_body=json_body,
                expected_status=expected_status,
                expected_body_subset=expect_subset,
            )
        )
        return resp

    def write_pact(self, path: Path) -> Path:
        path.parent.mkdir(parents=True, exist_ok=True)
        document = {
            "consumer": {"name": self.consumer},
            "provider": {"name": self.provider},
            "interactions": [i.to_json() for i in self.interactions],
        }
        path.write_text(json.dumps(document, indent=2) + "\n", encoding="utf-8")
        return path


@dataclass
class VerificationReport:
    """Outcome of replaying every pact interaction against the provider."""

    results: list[tuple[str, bool, list[str]]] = field(default_factory=list)

    @property
    def ok(self) -> bool:
        return all(passed for _, passed, _ in self.results)

    def failures(self) -> list[str]:
        return [f"{desc}: {problems}" for desc, passed, problems in self.results if not passed]

    def render(self) -> str:
        lines = []
        for desc, passed, problems in self.results:
            mark = "PASS" if passed else "FAIL"
            lines.append(f"  [{mark}] {desc}")
            for problem in problems:
                lines.append(f"         {problem}")
        return "\n".join(lines)


class ProviderVerifier:
    """Provider side: replay a pact file against a fresh application instance."""

    def __init__(
        self,
        app_factory: AppFactory,
        pact_path: Path,
        states: dict[str, StateHandler],
    ) -> None:
        self._app_factory = app_factory
        self._pact_path = pact_path
        self._states = states

    def verify(self) -> VerificationReport:
        document = json.loads(self._pact_path.read_text(encoding="utf-8"))
        report = VerificationReport()
        for raw in document["interactions"]:
            interaction = Interaction.from_json(raw)
            problems: list[str] = []
            app = self._app_factory()
            if interaction.provider_state is not None:
                handler = self._states.get(interaction.provider_state)
                if handler is None:
                    problems.append(f"unknown provider state: {interaction.provider_state!r}")
                else:
                    handler(app)
            if not problems:
                client = make_sync_client(app, base_url="http://provider.test")
                resp = client.request(
                    interaction.method,
                    interaction.path,
                    json=interaction.request_body,
                )
                if resp.status_code != interaction.expected_status:
                    problems.append(
                        f"status: expected {interaction.expected_status}, "
                        f"got {resp.status_code}"
                    )
                if interaction.expected_body_subset:
                    try:
                        body = resp.json()
                    except json.JSONDecodeError:
                        problems.append("response body is not JSON")
                    else:
                        problems.extend(subset_mismatches(
                            interaction.expected_body_subset, body
                        ))
            report.results.append((interaction.description, not problems, problems))
        return report


# --------------------------------------------------------------------------
# Demo consumer: the Northwind Robotics fulfillment service.
# --------------------------------------------------------------------------

def _seed_robot_7(app: Any) -> None:
    """Provider state: 'a robot with id 7 exists'."""
    app.state.db[7] = {
        "id": 7,
        "sku": "NWR-ARM-6",
        "name": "A6 Manipulator Arm",
        "price_cents": 129900,
    }


PROVIDER_STATES: dict[str, StateHandler] = {
    "a robot with id 7 exists": _seed_robot_7,
}


def record_demo_pact(pact_dir: Path) -> Path:
    """Run the fulfillment service's expectations and freeze them as a pact."""
    app = create_app()
    _seed_robot_7(app)  # consumer designs the 'GET' interaction against known data
    client = make_sync_client(app, base_url="http://consumer.test")
    recorder = PactRecorder(
        client, consumer="fulfillment-service", provider="robot-registry"
    )
    recorder.record(
        "create a robot",
        "POST",
        "/robots",
        json_body={"sku": "NWR-PLT-2", "name": "Pallet Mover 2", "price_cents": 540000},
        expected_status=201,
        expect_subset={"sku": "NWR-PLT-2", "name": "Pallet Mover 2", "price_cents": 540000},
    )
    recorder.record(
        "fetch robot 7",
        "GET",
        "/robots/7",
        provider_state="a robot with id 7 exists",
        expected_status=200,
        expect_subset={"id": 7, "sku": "NWR-ARM-6", "name": "A6 Manipulator Arm"},
    )
    return recorder.write_pact(pact_dir / "fulfillment-service-robot-registry.json")
\end{minted}

The pact file itself (excerpt) is intentionally boring --- boring is what
crosses team boundaries well:

\begin{minted}{json}
{
  "consumer": {"name": "fulfillment-service"},
  "provider": {"name": "robot-registry"},
  "interactions": [
    {"description": "fetch robot 7",
     "providerState": "a robot with id 7 exists",
     "request": {"method": "GET", "path": "/robots/7", "body": null},
     "response": {"status": 200,
                  "bodySubset": {"id": 7, "sku": "NWR-ARM-6"}}}
  ]
}
\end{minted}

\subsection{Listing 19.3 --- Contract Tests, Including the Breaking Change}
\label{sec:ch19-listing-pact-tests}

Four tests prove the framework works \emph{and has teeth}: the pact file is
well-formed; the healthy provider satisfies it; renaming \texttt{sku} to
\texttt{part\_number} makes verification fail with a message naming the
missing field; and an unknown provider state fails loudly rather than
silently passing.

\begin{minted}{python}
"""Contract tests for the mini pact framework.

Run:  pytest test_pact_mini.py -q
"""
from __future__ import annotations

import json
from pathlib import Path

from pact_mini import PROVIDER_STATES, ProviderVerifier, record_demo_pact
from robot_registry import create_app


def test_recorded_pact_file_is_wellformed(tmp_path: Path) -> None:
    pact = record_demo_pact(tmp_path)
    document = json.loads(pact.read_text(encoding="utf-8"))
    assert document["consumer"]["name"] == "fulfillment-service"
    assert document["provider"]["name"] == "robot-registry"
    assert [i["description"] for i in document["interactions"]] == [
        "create a robot",
        "fetch robot 7",
    ]
    assert document["interactions"][1]["providerState"] == "a robot with id 7 exists"


def test_provider_satisfies_the_pact(tmp_path: Path) -> None:
    pact = record_demo_pact(tmp_path)
    report = ProviderVerifier(create_app, pact, PROVIDER_STATES).verify()
    print("\n" + report.render())
    assert report.ok, report.failures()


def test_breaking_field_rename_is_caught(tmp_path: Path) -> None:
    """The provider renames ``sku`` -> ``part_number``: verification must fail."""
    pact = record_demo_pact(tmp_path)
    report = ProviderVerifier(
        lambda: create_app(breaking=True), pact, PROVIDER_STATES
    ).verify()
    print("\n" + report.render())
    assert not report.ok
    rendered = report.render()
    assert "sku" in rendered  # the missing field is named in the report


def test_unknown_provider_state_fails_loudly(tmp_path: Path) -> None:
    pact = record_demo_pact(tmp_path)
    document = json.loads(pact.read_text(encoding="utf-8"))
    document["interactions"][1]["providerState"] = "a robot with id 99 exists"
    pact.write_text(json.dumps(document), encoding="utf-8")
    report = ProviderVerifier(create_app, pact, PROVIDER_STATES).verify()
    assert not report.ok
    assert any("unknown provider state" in f for f in report.failures())
\end{minted}

\subsection{Listing 19.4 --- Property-Based Suite with Hypothesis}
\label{sec:ch19-listing-pbt}

Listing 19.4 encodes the three universal API invariants from
Section~\ref{sec:ch19-pbt}. The payload strategy is deliberately hostile:
wrong types, \texttt{None}, out-of-range integers, over- and under-length
strings. Whatever Hypothesis generates, the service must never 5xx, must
return schema-valid JSON on 201, and must return an RFC~9457 body on 422.
\texttt{derandomize=True} replays the identical examples on every run ---
determinism is a feature, not a limitation, in CI. The rule-based
\texttt{RobotRegistryMachine} then generates random create/read
\emph{sequences} and checks the API against a local model after every
step, catching sequence-dependent bugs single-call properties cannot see.

\begin{minted}{python}
"""Property-based tests for the robot registry, plus a stateful rule suite.

Invariants under test:

* ``POST /robots`` never returns 5xx, whatever the payload.
* A 201 response always validates against the published response schema.
* A 4xx response is always RFC 9457 problem-shaped.
* ``GET /robots/{id}`` never returns 5xx for arbitrary integer ids.
* Rule-based sequences (create -> read) stay consistent with a local model.

Deterministic: ``derandomize=True`` replays the same examples every run and
no Hypothesis example database is consulted.

Run:  pytest test_pbt_robots.py -q
"""
from __future__ import annotations

from typing import Any

import jsonschema
from hypothesis import given, settings, strategies as st
from hypothesis.stateful import (
    Bundle,
    RuleBasedStateMachine,
    invariant,
    rule,
)

from robot_registry import create_app, make_sync_client

CLIENT = make_sync_client(create_app(), base_url="http://pbt.test")

ROBOT_SCHEMA: dict[str, Any] = {
    "type": "object",
    "required": ["id", "sku", "name", "price_cents"],
    "properties": {
        "id": {"type": "integer", "minimum": 1},
        "sku": {"type": "string", "minLength": 3, "maxLength": 32},
        "name": {"type": "string", "minLength": 1, "maxLength": 80},
        "price_cents": {"type": "integer", "minimum": 0, "maximum": 10_000_000},
    },
    "additionalProperties": False,
}

PROBLEM_MEMBERS = {"type", "title", "status", "detail", "instance"}

# Deliberately hostile: wrong types, None, out-of-range ints, over/under strings.
HOSTILE_PAYLOADS = st.fixed_dictionaries(
    {
        "sku": st.one_of(st.text(min_size=0, max_size=40), st.integers(), st.none()),
        "name": st.one_of(st.text(max_size=100), st.booleans(), st.none()),
        "price_cents": st.one_of(
            st.integers(min_value=-1_000_000, max_value=20_000_000),
            st.text(max_size=10),
            st.none(),
        ),
    }
)

VALID_PAYLOADS = st.fixed_dictionaries(
    {
        "sku": st.text(
            alphabet=st.characters(categories=("Lu", "Ll", "Nd"), include_characters="-"),
            min_size=3,
            max_size=32,
        ),
        "name": st.text(
            alphabet=st.characters(categories=("Lu", "Ll", "Nd", "Zs"), include_characters="-"),
            min_size=1,
            max_size=80,
        ),
        "price_cents": st.integers(min_value=0, max_value=10_000_000),
    }
)

DETERMINISTIC = dict(derandomize=True, database=None, deadline=None)


def assert_problem_shape(body: Any, status: int) -> None:
    assert isinstance(body, dict), f"4xx body must be a JSON object, got {body!r}"
    missing = PROBLEM_MEMBERS - body.keys()
    assert not missing, f"problem details missing members: {missing}"
    assert body["status"] == status


@given(payload=HOSTILE_PAYLOADS)
@settings(max_examples=60, **DETERMINISTIC)
def test_create_robot_never_5xx_and_always_schema_conformant(payload: dict) -> None:
    resp = CLIENT.post("/robots", json=payload)
    assert resp.status_code < 500, f"5xx for payload {payload!r}: {resp.text}"
    if resp.status_code == 201:
        jsonschema.validate(resp.json(), ROBOT_SCHEMA)
    else:
        assert resp.status_code == 422, f"unexpected status {resp.status_code}"
        assert resp.headers["content-type"].startswith("application/problem+json")
        assert_problem_shape(resp.json(), 422)


@given(robot_id=st.integers(min_value=-1_000, max_value=1_000))
@settings(max_examples=50, **DETERMINISTIC)
def test_get_robot_never_5xx(robot_id: int) -> None:
    resp = CLIENT.get(f"/robots/{robot_id}")
    assert resp.status_code in (200, 404), f"unexpected status {resp.status_code}"
    if resp.status_code == 200:
        jsonschema.validate(resp.json(), ROBOT_SCHEMA)
    else:
        assert_problem_shape(resp.json(), 404)


class RobotRegistryMachine(RuleBasedStateMachine):
    """Rule-based stateful test: random create/read sequences checked against
    a local model of the registry."""

    created_ids: Bundle[int] = Bundle("created_ids")

    def __init__(self) -> None:
        super().__init__()
        self.client = make_sync_client(create_app(), base_url="http://pbt-stateful.test")
        self.model: dict[int, dict[str, Any]] = {}

    @rule(target=created_ids, payload=VALID_PAYLOADS)
    def create_robot(self, payload: dict[str, Any]) -> int:
        resp = self.client.post("/robots", json=payload)
        assert resp.status_code == 201, resp.text
        body = resp.json()
        jsonschema.validate(body, ROBOT_SCHEMA)
        self.model[body["id"]] = body
        return body["id"]

    @rule(robot_id=created_ids)
    def read_back_matches_model(self, robot_id: int) -> None:
        resp = self.client.get(f"/robots/{robot_id}")
        assert resp.status_code == 200
        assert resp.json() == self.model[robot_id]

    @invariant()
    def every_created_robot_is_readable(self) -> None:
        for robot_id in self.model:
            resp = self.client.get(f"/robots/{robot_id}")
            assert resp.status_code == 200, f"robot {robot_id} vanished"


RobotRegistryMachine.TestCase.settings = settings(
    max_examples=15, stateful_step_count=20, **DETERMINISTIC
)
TestRobotRegistry = RobotRegistryMachine.TestCase
\end{minted}

\subsection{Listing 19.5 --- A Stdlib OpenAPI Fuzzer with Shrinking}
\label{sec:ch19-listing-fuzz}

Where Hypothesis owns the strategy machinery, Listing 19.5
exposes it: from a request-body schema (a stand-in for an excerpt of the
published OpenAPI document~\cite{openapi31}) it derives boundary and
malformed candidates per property --- \texttt{minLength}$-1$,
\texttt{maxLength}$+1$, \texttt{minimum}$-1$, wrong types, \texttt{None},
pathological lengths --- plus missing-required, extra-property, and
non-object-body cases. Every response is classified: 5xx is always a
failure, 4xx must be problem-shaped, 2xx must validate. The greedy
\texttt{shrink} first tries dropping keys, then minimizes values to the
smallest schema-valid stand-ins --- the second move matters, because
dropping a \emph{required} key flips a 500 into a 422 and stalls a
naive shrinker.

\begin{minted}{python}
"""Stdlib-only OpenAPI-driven fuzzer: Schemathesis in miniature.

Given the request-body schema of ``POST /robots`` (a stand-in for an excerpt
of the published OpenAPI document), generate boundary and malformed
payloads, fire them at the app in-process, and classify every response:

* 2xx  -> must validate against the response schema;
* 4xx  -> must be RFC 9457 problem-shaped;
* 5xx  -> always a failure.

Includes a greedy shrinker: when a failing payload is found, strip it down
to a minimal reproducer so the report is actionable.

Run the demo:  python fuzz_openapi_mini.py
"""
from __future__ import annotations

import copy
import json
from dataclasses import dataclass
from typing import Any, Callable

import httpx
import jsonschema

from robot_registry import create_app, make_sync_client

# Excerpt of the published OpenAPI 3.1 document: POST /robots request body.
REQUEST_SCHEMA: dict[str, Any] = {
    "type": "object",
    "required": ["sku", "name", "price_cents"],
    "properties": {
        "sku": {"type": "string", "minLength": 3, "maxLength": 32},
        "name": {"type": "string", "minLength": 1, "maxLength": 80},
        "price_cents": {"type": "integer", "minimum": 0, "maximum": 10_000_000},
    },
    "additionalProperties": False,
}

RESPONSE_SCHEMA: dict[str, Any] = {
    "type": "object",
    "required": ["id", "sku", "name", "price_cents"],
    "properties": {
        "id": {"type": "integer", "minimum": 1},
        "sku": {"type": "string", "minLength": 3, "maxLength": 32},
        "name": {"type": "string", "minLength": 1, "maxLength": 80},
        "price_cents": {"type": "integer", "minimum": 0, "maximum": 10_000_000},
    },
    "additionalProperties": False,
}

PROBLEM_MEMBERS = {"type", "title", "status", "detail", "instance"}
VALID_BASE: dict[str, Any] = {"sku": "NWR-ARM-6", "name": "A6 Manipulator Arm", "price_cents": 129900}


def boundary_candidates(prop_schema: dict[str, Any]) -> list[Any]:
    """Boundary and malformed candidates for one property schema."""
    ptype = prop_schema.get("type")
    if ptype == "string":
        lo = prop_schema.get("minLength", 0)
        hi = prop_schema.get("maxLength", 16)
        candidates = ["A" * lo, "A" * hi, "", None, 123, True]
        if lo > 0:
            candidates.append("A" * (lo - 1))
        candidates.append("A" * (hi + 1))
        candidates.append("Z" * 512)  # pathological length
        return candidates
    if ptype == "integer":
        lo = prop_schema.get("minimum", 0)
        hi = prop_schema.get("maximum", 100)
        return [lo, hi, lo - 1, hi + 1, 0, -1, 666, 2**31, "10", None, True, 1.5]
    return [None]


def fuzz_cases(schema: dict[str, Any]) -> list[tuple[str, Any]]:
    """All (label, payload) cases derived from the request schema."""
    cases: list[tuple[str, Any]] = []
    properties = schema["properties"]
    for name, prop_schema in properties.items():
        for value in boundary_candidates(prop_schema):
            payload = dict(VALID_BASE, **{name: value})
            cases.append((f"{name}={value!r}"[:60], payload))
    for name in schema.get("required", []):
        payload = {k: v for k, v in VALID_BASE.items() if k != name}
        cases.append((f"missing required {name!r}", payload))
    cases.append(("extra property", dict(VALID_BASE, legacy_field="x")))
    cases.append(("empty object", {}))
    cases.append(("body is a list", [1, 2, 3]))
    cases.append(("body is a string", "not-an-object"))
    cases.append(("body is null", None))
    return cases


@dataclass
class FuzzFailure:
    label: str
    payload: Any
    status: int
    reason: str


def classify(label: str, payload: Any, resp: httpx.Response) -> FuzzFailure | None:
    if resp.status_code >= 500:
        return FuzzFailure(label, payload, resp.status_code, "server error (5xx)")
    if resp.status_code >= 400:
        body = resp.json()
        missing = PROBLEM_MEMBERS - set(body)
        if missing or body.get("status") != resp.status_code:
            return FuzzFailure(
                label, payload, resp.status_code, f"4xx not RFC 9457 shaped: {body!r}"
            )
        return None
    try:
        jsonschema.validate(resp.json(), RESPONSE_SCHEMA)
    except jsonschema.ValidationError as exc:
        return FuzzFailure(label, payload, resp.status_code, f"2xx violates schema: {exc.message}")
    return None


def run_fuzz(
    client: httpx.Client, cases: list[tuple[str, Any]]
) -> tuple[int, list[FuzzFailure]]:
    failures: list[FuzzFailure] = []
    for label, payload in cases:
        resp = client.post("/robots", json=payload)
        failure = classify(label, payload, resp)
        if failure is not None:
            failures.append(failure)
    return len(cases), failures


def minimal_values(prop_schema: dict[str, Any]) -> list[Any]:
    """Smallest schema-valid stand-ins for a property (value minimization)."""
    ptype = prop_schema.get("type")
    if ptype == "string":
        return ["A" * prop_schema.get("minLength", 0)]
    if ptype == "integer":
        return [prop_schema.get("minimum", 0)]
    return [None]


def shrink(
    client: httpx.Client,
    payload: dict[str, Any],
    is_failure: Callable[[httpx.Response], bool],
    schema: dict[str, Any] | None = None,
) -> dict[str, Any]:
    """Greedy minimization: drop keys, then minimize values, while the
    failure still reproduces. Dropping a *required* key usually turns a 500
    into a 422, so value minimization is what isolates the trigger field."""
    minimal = copy.deepcopy(payload)
    for key in sorted(minimal):
        candidate = {k: v for k, v in minimal.items() if k != key}
        if is_failure(client.post("/robots", json=candidate)):
            minimal = candidate
    if schema is not None:
        for key, prop_schema in schema.get("properties", {}).items():
            if key not in minimal:
                continue
            for candidate_value in minimal_values(prop_schema):
                trial = dict(minimal, **{key: candidate_value})
                if is_failure(client.post("/robots", json=trial)):
                    minimal = trial
                    break
    return minimal


def demo() -> None:
    cases = fuzz_cases(REQUEST_SCHEMA)
    healthy = make_sync_client(create_app(), base_url="http://fuzz.test")
    total, failures = run_fuzz(healthy, cases)
    print(f"healthy app:    {total} cases, {len(failures)} failures")

    buggy = make_sync_client(create_app(bug_at_price=666), base_url="http://fuzz.test")
    total, failures = run_fuzz(buggy, cases)
    print(f"buggy app:      {total} cases, {len(failures)} failures")
    if failures:
        first = failures[0]
        minimal = shrink(
            buggy, first.payload, lambda r: r.status_code >= 500, REQUEST_SCHEMA
        )
        print(f"first failure:  {first.label} -> {first.status} ({first.reason})")
        print(f"shrunk to:      {json.dumps(minimal, sort_keys=True)}")


if __name__ == "__main__":
    demo()
\end{minted}

Its tests (the companion test file) pin the report-hygiene rules of
Section~\ref{sec:ch19-pitfalls}: prove the fuzzer actually exercised the
validation boundary (count the 4xxs), prove it finds a planted bug, and
prove the shrinker converges to the minimal reproducer.

\begin{minted}{python}
"""Tests for the mini OpenAPI fuzzer.

Run:  pytest test_fuzz_openapi_mini.py -q
"""
from __future__ import annotations

import httpx

from fuzz_openapi_mini import REQUEST_SCHEMA, fuzz_cases, run_fuzz, shrink
from robot_registry import create_app, make_sync_client


def _client(**app_kwargs: object) -> httpx.Client:
    return make_sync_client(
        create_app(**app_kwargs),  # type: ignore[arg-type]
        base_url="http://fuzz.test",
    )


def test_fuzzer_finds_no_failures_on_healthy_app() -> None:
    cases = fuzz_cases(REQUEST_SCHEMA)
    total, failures = run_fuzz(_client(), cases)
    assert total == len(cases) and total > 30  # report hygiene: prove coverage
    assert failures == [], [f"{f.label}: {f.reason}" for f in failures]


def test_fuzzer_finds_and_shrinks_the_planted_bug() -> None:
    cases = fuzz_cases(REQUEST_SCHEMA)
    client = _client(bug_at_price=666)
    _, failures = run_fuzz(client, cases)
    assert failures, "fuzzer must find the planted 500"
    assert all(f.status == 500 for f in failures)
    first = failures[0]
    minimal = shrink(client, first.payload, lambda r: r.status_code >= 500, REQUEST_SCHEMA)
    # Required keys cannot be dropped (that flips 500 -> 422), so the minimal
    # reproducer isolates the trigger field with minimal valid stand-ins:
    assert minimal == {"sku": "AAA", "name": "A", "price_cents": 666}


def test_every_4xx_is_problem_shaped() -> None:
    """A focused regression of the RFC 9457 invariant on the healthy app."""
    cases = fuzz_cases(REQUEST_SCHEMA)
    client = _client()
    saw_4xx = 0
    for label, payload in cases:
        resp = client.post("/robots", json=payload)
        if 400 <= resp.status_code < 500:
            saw_4xx += 1
            body = resp.json()
            assert body["status"] == resp.status_code, label
            assert resp.headers["content-type"].startswith("application/problem+json")
    assert saw_4xx > 10  # the fuzzer really exercised the validation boundary
\end{minted}

\subsection{Listing 19.6 --- An asyncio Load Harness with Knee-Point Sweep}
\label{sec:ch19-listing-load}

Listing 19.6 implements both workload models against the
in-process app: \texttt{open\_loop} schedules sends at fixed virtual
offsets (offered load is independent of completions --- no coordinated
omission), while \texttt{closed\_loop} runs a fixed worker pool. A
fixed-bucket histogram estimates percentiles without storing every sample
(the same machinery as Chapter~\ref{ch:observability}). The demo workload
hides simulated 20~ms of async work behind a semaphore of four slots, so
its saturation point is real and findable: capacity $\approx c/S =
4/0.02 = 200$ req/s, exactly where the printed knee appears.
\texttt{littles\_law\_check} closes the loop with
Theorem~\ref{thm:little}: with eight closed-model workers it measures
$L_{\text{est}} = \lambda \bar{W} \approx 8$. The sweep prints a table;
nothing is asserted about wall-clock numbers.

\begin{minted}{python}
"""Stdlib + asyncio load harness for offline API load testing.

Implements both classic workload models against an in-process ASGI app:

* open model   -- requests are scheduled at a fixed rate (requests/sec),
                  independent of how fast responses come back;
* closed model -- a fixed number of workers each run request-response
                  loops (think: a fixed pool of users).

Includes a bucketed latency histogram with percentile estimation, a
saturation sweep that prints a knee-point table, and a Little's Law
cross-check (L = lambda * W) for the closed model.

Latency numbers are *illustrative* (they depend on your machine); the
harness never asserts on wall-clock values.

Run the demo:  python load_harness.py
"""
from __future__ import annotations

import asyncio
import bisect
import time
from dataclasses import dataclass, field
from typing import Any

import httpx
from fastapi import FastAPI

BOUNDS_MS: tuple[int, ...] = (1, 2, 5, 10, 25, 50, 100, 250, 500, 1_000, 2_500, 5_000)


@dataclass
class Histogram:
    """Fixed-bucket latency histogram; percentiles read the bucket ceiling."""

    bounds: tuple[int, ...] = BOUNDS_MS
    counts: list[int] = field(init=False)
    errors: int = 0
    total_ms: float = 0.0

    def __post_init__(self) -> None:
        self.counts = [0] * (len(self.bounds) + 1)

    def record_ms(self, ms: float) -> None:
        self.counts[bisect.bisect_right(self.bounds, ms)] += 1
        self.total_ms += ms

    @property
    def total(self) -> int:
        return sum(self.counts)

    def percentile_ms(self, p: float) -> float:
        if self.total == 0:
            return 0.0
        rank = max(1, int(p / 100.0 * self.total + 0.999999))
        seen = 0
        for i, count in enumerate(self.counts):
            seen += count
            if seen >= rank:
                return float(self.bounds[i]) if i < len(self.bounds) else float("inf")
        return float("inf")

    def mean_ms(self) -> float:
        return self.total_ms / self.total if self.total else 0.0


async def _send(client: httpx.AsyncClient, hist: Histogram, path: str) -> None:
    started = time.perf_counter()
    try:
        resp = await client.get(path)
        await resp.aread()
        if resp.status_code >= 500:
            hist.errors += 1
    except (httpx.HTTPError, OSError):
        hist.errors += 1
    finally:
        hist.record_ms((time.perf_counter() - started) * 1_000.0)


async def open_loop(
    client: httpx.AsyncClient, *, rps: float, duration_s: float, path: str = "/work"
) -> Histogram:
    """Open model: schedule sends at fixed virtual offsets from the start."""
    hist = Histogram()
    period = 1.0 / rps
    start = time.perf_counter()
    tasks: list[asyncio.Task[None]] = []
    k = 0
    while k * period < duration_s:  # deterministic scheduled count
        scheduled = start + k * period
        delay = scheduled - time.perf_counter()
        if delay > 0:
            await asyncio.sleep(delay)
        tasks.append(asyncio.create_task(_send(client, hist, path)))
        k += 1
    if tasks:
        await asyncio.gather(*tasks)
    return hist


async def closed_loop(
    client: httpx.AsyncClient, *, workers: int, duration_s: float, path: str = "/work"
) -> Histogram:
    """Closed model: each worker issues the next request only after the
    previous response arrives."""
    hist = Histogram()
    deadline = time.perf_counter() + duration_s

    async def worker() -> None:
        while time.perf_counter() < deadline:
            await _send(client, hist, path)

    await asyncio.gather(*(worker() for _ in range(workers)))
    return hist


def create_workload_app(*, work_ms: int = 20, max_concurrent: int = 4) -> FastAPI:
    """Demo workload: simulated async work behind a concurrency semaphore,
    so the service has a real, findable saturation point."""
    app = FastAPI(title="Load harness demo workload")
    semaphore = asyncio.Semaphore(max_concurrent)

    @app.get("/work")
    async def work() -> dict[str, str]:
        async with semaphore:
            await asyncio.sleep(work_ms / 1_000.0)
        return {"status": "ok"}

    return app


@dataclass
class SweepRow:
    rps: int
    completed: int
    p50_ms: float
    p99_ms: float
    error_rate: float
    is_knee: bool


async def saturation_sweep(
    app: Any,
    levels: tuple[int, ...] = (25, 50, 100, 150, 200, 300, 400),
    duration_s: float = 0.3,
) -> list[SweepRow]:
    """Ramp offered RPS; flag the knee: first level where p99 explodes
    (10x baseline) or errors appear."""
    rows: list[SweepRow] = []
    baseline_p50: float | None = None
    async with httpx.AsyncClient(
        transport=httpx.ASGITransport(app=app), base_url="http://load.test"
    ) as client:
        for rps in levels:
            hist = await open_loop(client, rps=rps, duration_s=duration_s)
            p50 = hist.percentile_ms(50)
            p99 = hist.percentile_ms(99)
            error_rate = hist.errors / hist.total if hist.total else 1.0
            if baseline_p50 is None:
                baseline_p50 = p50 if p50 > 0 else 1.0
            is_knee = (p99 > 10.0 * baseline_p50) or error_rate > 0.0
            rows.append(SweepRow(rps, hist.total, p50, p99, error_rate, is_knee))
    return rows


def render_sweep(rows: list[SweepRow]) -> str:
    lines = ["  rps | completed |  p50 ms |   p99 ms | err % | knee",
             "------+-----------+---------+----------+-------+-----"]
    for row in rows:
        lines.append(
            f"{row.rps:5} | {row.completed:9} | {row.p50_ms:7.0f} | "
            f"{row.p99_ms:8.0f} | {100 * row.error_rate:5.1f} | "
            f"{'<< KNEE' if row.is_knee else ''}"
        )
    return "\n".join(lines)


async def littles_law_check(workers: int = 8, duration_s: float = 0.5) -> str:
    """Closed model: verify L = lambda * W against the configured concurrency."""
    app = create_workload_app(work_ms=20, max_concurrent=10_000)  # unthrottled
    async with httpx.AsyncClient(
        transport=httpx.ASGITransport(app=app), base_url="http://load.test"
    ) as client:
        hist = await closed_loop(client, workers=workers, duration_s=duration_s)
    lam = hist.total / duration_s            # throughput lambda (requests/s)
    w = hist.mean_ms() / 1_000.0             # mean residence time W (s)
    l_est = lam * w
    return (
        f"closed model: workers={workers} throughput={lam:.0f} req/s "
        f"mean={hist.mean_ms():.1f} ms -> L = lambda*W = {l_est:.1f} "
        f"(configured concurrency: {workers})"
    )


async def _demo() -> None:
    print("saturation sweep (illustrative numbers, machine-dependent):")
    rows = await saturation_sweep(create_workload_app())
    print(render_sweep(rows))
    print(await littles_law_check())


if __name__ == "__main__":
    asyncio.run(_demo())
\end{minted}

Illustrative output (your numbers will differ; the shape will not):

\begin{minted}{text}
  rps | completed |  p50 ms |   p99 ms | err % | knee
------+-----------+---------+----------+-------+-----
   25 |         8 |      50 |       50 |   0.0 |
  100 |        30 |      50 |       50 |   0.0 |
  200 |        60 |     250 |      250 |   0.0 |
  400 |       120 |     500 |     1000 |   0.0 | << KNEE
closed model: workers=8 throughput=270 req/s mean=30.5 ms
              -> L = lambda*W = 8.2 (configured concurrency: 8)
\end{minted}

\subsection{Listing 19.7 --- Fault-Injection Middleware and Resilient Client}
\label{sec:ch19-listing-faults}

Listing 19.7 is the chaos toolkit. \texttt{FaultScript} maps
request index to fault --- deterministic by construction, no dice. The pure
ASGI \texttt{FaultInjectionMiddleware} implements three fault classes from
the taxonomy: \texttt{LATENCY} (sleep on the \emph{virtual} clock, then
serve), \texttt{ERROR} (answer directly with an RFC~9457-shaped 503),
\texttt{ABORT} (raise before any response, simulating a reset connection).
The client under test, \texttt{ResilientTransport}, is a real
\texttt{httpx} transport: bounded retries with exponential backoff on the
injected sleeper, and a circuit breaker that fails fast after
\texttt{breaker\_threshold} consecutive exhausted calls.

\begin{minted}{python}
"""Deterministic fault-injection middleware and a resilient client transport.

Fault classes (see the chapter's fault taxonomy):

* ``LATENCY`` -- sleep (on a *virtual* clock) before calling the app;
* ``ERROR``   -- answer directly with an RFC 9457-shaped 5xx;
* ``ABORT``   -- drop the connection (raise before any response).

The :class:`VirtualClock` sleeper never waits in real time: tests exercise
multi-second retry backoffs in microseconds and assert on *recorded*
virtual delays, never on the wall clock.

The :class:`ResilientTransport` is the client under test: bounded retries
with exponential backoff plus a circuit breaker that fails fast after
``breaker_threshold`` consecutive exhausted calls.
"""
from __future__ import annotations

from dataclasses import dataclass
from enum import StrEnum
from typing import Any, Awaitable, Callable, Sequence

import httpx

Sleeper = Callable[[float], Awaitable[None]]


class FaultKind(StrEnum):
    LATENCY = "latency"
    ERROR = "error"
    ABORT = "abort"


@dataclass(frozen=True)
class Fault:
    kind: FaultKind
    delay_ms: int = 0
    status: int = 503


class FaultScript:
    """Deterministic fault schedule: ``fault_for(i)`` depends only on the
    request index ``i``; beyond the script the service is healthy."""

    def __init__(self, entries: Sequence[Fault | None]) -> None:
        self._entries = list(entries)

    def fault_for(self, index: int) -> Fault | None:
        return self._entries[index] if index < len(self._entries) else None


class FaultInjectionMiddleware:
    """Pure ASGI middleware applying the scripted fault per request."""

    def __init__(self, app: Any, script: FaultScript, sleeper: Sleeper) -> None:
        self._app = app
        self._script = script
        self._sleeper = sleeper
        self.requests_seen = 0

    async def __call__(self, scope: dict, receive: Callable, send: Callable) -> None:
        if scope["type"] != "http":
            await self._app(scope, receive, send)
            return
        fault = self._script.fault_for(self.requests_seen)
        self.requests_seen += 1
        if fault is None or fault.kind is FaultKind.LATENCY:
            if fault is not None:
                await self._sleeper(fault.delay_ms / 1_000.0)
            await self._app(scope, receive, send)
            return
        if fault.kind is FaultKind.ERROR:
            body = (
                b'{"type":"https://northwind-robotics.example/problems/injected",'
                b'"title":"Injected failure","status":%d,'
                b'"detail":"chaos middleware","instance":"%s"}'
                % (fault.status, scope.get("path", "/").encode())
            )
            await send({
                "type": "http.response.start",
                "status": fault.status,
                "headers": [(b"content-type", b"application/problem+json")],
            })
            await send({"type": "http.response.body", "body": body})
            return
        raise ConnectionResetError("injected abort")  # FaultKind.ABORT


class VirtualClock:
    """Sleeper that never waits: records delays, advances virtual time."""

    def __init__(self) -> None:
        self.now = 0.0
        self.sleeps: list[float] = []

    async def sleep(self, seconds: float) -> None:
        self.sleeps.append(seconds)
        self.now += seconds


class BreakerOpenError(httpx.TransportError):
    """Raised when the circuit breaker short-circuits a call."""


@dataclass(frozen=True)
class RetryPolicy:
    retries: int = 3
    base_backoff_s: float = 0.1
    breaker_threshold: int = 3


class ResilientTransport(httpx.AsyncBaseTransport):
    """Bounded retries with virtual-clock backoff + circuit breaker."""

    def __init__(
        self,
        inner: httpx.AsyncBaseTransport,
        policy: RetryPolicy | None = None,
        sleeper: Sleeper | None = None,
    ) -> None:
        self._inner = inner
        self._policy = policy or RetryPolicy()
        self._sleeper = sleeper or _real_sleep
        self._consecutive_failures = 0

    async def handle_async_request(self, request: httpx.Request) -> httpx.Response:
        if self._consecutive_failures >= self._policy.breaker_threshold:
            raise BreakerOpenError("circuit breaker open: failing fast")
        policy = self._policy
        last_resp: httpx.Response | None = None
        last_exc: Exception | None = None
        for attempt in range(policy.retries + 1):
            try:
                resp = await self._inner.handle_async_request(request)
            except (httpx.TransportError, OSError) as exc:
                last_exc, resp = exc, None  # ABORT class: no response at all
            if resp is not None and resp.status_code < 500:
                self._consecutive_failures = 0
                return resp
            if resp is not None:
                await resp.aread()
                last_resp = resp
            if attempt < policy.retries:
                await self._sleeper(policy.base_backoff_s * (2**attempt))
        self._consecutive_failures += 1
        if last_resp is not None:
            return last_resp
        assert last_exc is not None
        raise last_exc


async def _real_sleep(seconds: float) -> None:
    import asyncio

    await asyncio.sleep(seconds)
\end{minted}

The tests (the companion test file) run multi-second backoff schedules
in microseconds and assert only on recorded virtual time and call counts:
recovery from two injected 5xxs with backoff \texttt{[0.1, 0.2]}, a 250~ms
latency fault that never touches the wall clock, abort storms surfacing as
transport errors, a breaker that opens after two exhausted calls and
short-circuits the third (proven by the middleware's request counter), and
a healthy passthrough that burns zero retries.

\begin{minted}{python}
"""Deterministic resilience tests: every fault class, zero wall-clock waits.

Run:  pytest test_fault_injection.py -q
"""
from __future__ import annotations

import asyncio

import httpx
import pytest

from fault_injection import (
    BreakerOpenError,
    Fault,
    FaultInjectionMiddleware,
    FaultKind,
    FaultScript,
    ResilientTransport,
    RetryPolicy,
    VirtualClock,
)
from robot_registry import create_app

SEEDED_ROBOT = {"id": 1, "sku": "NWR-ARM-6", "name": "A6 Manipulator Arm", "price_cents": 129900}


def make_rig(
    script: FaultScript, policy: RetryPolicy | None = None
) -> tuple[httpx.AsyncClient, VirtualClock, FaultInjectionMiddleware]:
    clock = VirtualClock()
    app = create_app()
    app.state.db[1] = dict(SEEDED_ROBOT)
    middleware = FaultInjectionMiddleware(app, script, clock.sleep)
    inner = httpx.ASGITransport(app=middleware)
    transport = ResilientTransport(inner, policy or RetryPolicy(), clock.sleep)
    client = httpx.AsyncClient(transport=transport, base_url="http://faults.test")
    return client, clock, middleware


def run(coro):  # keep pytest asyncio-plugin free
    return asyncio.run(coro)


def test_recovers_from_transient_errors() -> None:
    async def scenario() -> tuple[int, list[float]]:
        client, clock, _ = make_rig(FaultScript([Fault(FaultKind.ERROR), Fault(FaultKind.ERROR), None]))
        async with client:
            resp = await client.get("/robots/1")
        return resp.status_code, clock.sleeps

    status, sleeps = run(scenario())
    assert status == 200
    assert sleeps == [0.1, 0.2]  # exponential backoff on the virtual clock


def test_latency_fault_is_injected_on_virtual_time_only() -> None:
    async def scenario() -> tuple[int, float]:
        client, clock, _ = make_rig(
            FaultScript([Fault(FaultKind.LATENCY, delay_ms=250)]),
            RetryPolicy(retries=0),
        )
        async with client:
            resp = await client.get("/robots/1")
        return resp.status_code, clock.now

    status, virtual_now = run(scenario())
    assert status == 200
    assert virtual_now == 0.25  # recorded, not slept


def test_abort_retries_then_surfaces_transport_error() -> None:
    async def scenario() -> list[float]:
        client, clock, _ = make_rig(
            FaultScript([Fault(FaultKind.ABORT), Fault(FaultKind.ABORT)]),
            RetryPolicy(retries=1),
        )
        async with client:
            with pytest.raises(ConnectionResetError):
                await client.get("/robots/1")
        return clock.sleeps

    assert run(scenario()) == [0.1]


def test_circuit_breaker_fails_fast_after_threshold() -> None:
    async def scenario() -> tuple[int, int]:
        client, _, middleware = make_rig(
            FaultScript([Fault(FaultKind.ERROR)] * 5),
            RetryPolicy(retries=0, breaker_threshold=2),
        )
        async with client:
            first = await client.get("/robots/1")   # 503 -> failures = 1
            second = await client.get("/robots/1")  # 503 -> failures = 2 (threshold)
            assert (first.status_code, second.status_code) == (503, 503)
            with pytest.raises(BreakerOpenError):
                await client.get("/robots/1")       # short-circuited
        return middleware.requests_seen, client is not None and 1

    requests_seen, _ = run(scenario())
    assert requests_seen == 2  # third call never reached the middleware


def test_healthy_passthrough_needs_no_retries() -> None:
    async def scenario() -> tuple[int, list[float], int]:
        client, clock, middleware = make_rig(FaultScript([]))
        async with client:
            resp = await client.get("/robots/1")
        return resp.status_code, clock.sleeps, middleware.requests_seen

    status, sleeps, seen = run(scenario())
    assert (status, sleeps, seen) == (200, [], 1)
\end{minted}

\subsection{Listing 19.8 --- The Quality Gate}
\label{sec:ch19-listing-gate}

\texttt{ci\_gate.py} composes the chapter into one blocking check. Each
sub-gate is \emph{self-checking}: the contract gate additionally verifies
that the deliberately broken provider variant is rejected, and the fuzz
gate requires that the planted-bug variant produces failures --- a gate
that cannot fail is not a gate. The perf smoke asserts structure (all
scheduled requests completed, zero 5xx), never wall-clock latency. Exit
code 0 means merge; 1 means read the report.

\begin{minted}{python}
"""CI quality gate: contract verify + fuzz smoke + perf smoke -> exit code.

Each sub-gate returns (passed, detail). The contract and fuzz gates include
a *self-check*: they also run against a deliberately broken variant and
require that the gate catches it -- a gate that cannot fail is not a gate.

Run:  python ci_gate.py        (exit code 0 = PASS, 1 = FAIL)
"""
from __future__ import annotations

import asyncio
import sys
from pathlib import Path

import httpx

from fuzz_openapi_mini import REQUEST_SCHEMA, fuzz_cases, run_fuzz
from load_harness import create_workload_app, open_loop
from pact_mini import PROVIDER_STATES, ProviderVerifier, record_demo_pact
from robot_registry import create_app, make_sync_client

PACT_DIR = Path(__file__).resolve().parent / "pacts"


def contract_gate() -> tuple[bool, str]:
    pact = record_demo_pact(PACT_DIR)
    healthy = ProviderVerifier(create_app, pact, PROVIDER_STATES).verify()
    broken = ProviderVerifier(lambda: create_app(breaking=True), pact, PROVIDER_STATES).verify()
    detail = (
        f"{len(healthy.results)} interactions, {len(healthy.failures())} mismatches; "
        f"self-check: breaking variant rejected with "
        f"{len(broken.failures())} mismatch(es)"
    )
    return healthy.ok and not broken.ok, detail


def fuzz_gate() -> tuple[bool, str]:
    cases = fuzz_cases(REQUEST_SCHEMA)
    healthy_client = make_sync_client(create_app(), base_url="http://gate.test")
    total, failures = run_fuzz(healthy_client, cases)
    buggy_client = make_sync_client(create_app(bug_at_price=666), base_url="http://gate.test")
    _, self_check_failures = run_fuzz(buggy_client, cases)
    detail = (
        f"{total} cases, {len(failures)} failures; "
        f"self-check: planted bug found ({len(self_check_failures)} failing case(s))"
    )
    return not failures and bool(self_check_failures), detail


def perf_gate() -> tuple[bool, str]:
    """Perf smoke: short open-model burst. Asserts *structure* (all scheduled
    requests completed, zero 5xx), never wall-clock latency thresholds."""

    async def burst() -> tuple[int, int]:
        app = create_workload_app()
        async with httpx.AsyncClient(
            transport=httpx.ASGITransport(app=app), base_url="http://gate.test"
        ) as client:
            hist = await open_loop(client, rps=50, duration_s=0.4)
        return hist.total, hist.errors

    total, errors = asyncio.run(burst())
    expected = int(50 * 0.4)  # deterministic scheduled count modulo float edge
    ok = errors == 0 and abs(total - expected) <= 1
    return ok, f"{total} requests (~{expected} scheduled), {errors} server errors"


GATES = (
    ("contract verify", contract_gate),
    ("fuzz smoke", fuzz_gate),
    ("perf smoke", perf_gate),
)


def main() -> int:
    print("== API quality gate ==")
    all_ok = True
    for name, gate in GATES:
        try:
            passed, detail = gate()
        except Exception as exc:  # a crashing gate is a failing gate
            passed, detail = False, f"gate crashed: {exc!r}"
        all_ok &= passed
        print(f"[{'PASS' if passed else 'FAIL'}] {name:<16} {detail}")
    print(f"GATE RESULT: {'PASS' if all_ok else 'FAIL'}")
    return 0 if all_ok else 1


if __name__ == "__main__":
    sys.exit(main())
\end{minted}

\begin{minted}{text}
== API quality gate ==
[PASS] contract verify  2 interactions, 0 mismatches; self-check: breaking
                        variant rejected with 2 mismatch(es)
[PASS] fuzz smoke       38 cases, 0 failures; self-check: planted bug
                        found (1 failing case(s))
[PASS] perf smoke       20 requests (~20 scheduled), 0 server errors
GATE RESULT: PASS
\end{minted}

\subsection{Listing 19.9 --- The Pipeline as Code}
\label{sec:ch19-listing-yaml}

The GitHub Actions workflow is PowerShell-native (Windows runner,
\texttt{pwsh} shells, backslash paths --- no POSIX assumptions):

\begin{minted}{yaml}
name: api-quality-gates
on:
  pull_request:
  push:
    branches: [main]

jobs:
  quality:
    runs-on: windows-latest
    steps:
      - uses: actions/checkout@v4
      - uses: actions/setup-python@v5
        with:
          python-version: "3.11"
      - name: Install dependencies
        shell: pwsh
        run: |
          python -m venv .venv
          .\.venv\Scripts\Activate.ps1
          pip install -r code\ch19\requirements.txt
      - name: Unit, property, contract, fuzz, resilience tests
        shell: pwsh
        run: pytest code\ch19 -q
      - name: Blocking quality gate
        shell: pwsh
        run: python code\ch19\ci_gate.py
\end{minted}

%% ----------------------------------------------------------------------------
\section{Common Pitfalls \& Anti-Patterns}
\label{sec:ch19-pitfalls}

Each of these is a production scar, not a stylistic preference.

\begin{enumerate}[label=\textbf{P-\arabic*.}, leftmargin=1.9cm]
  \item \textbf{Testing implementation instead of behavior.} A test that
        asserts \texttt{validator() was called with X} breaks on every
        refactor and catches nothing. Assert on the wire contract: status,
        body shape, headers. The refactor-safety of a suite is its
        behavioral surface area.
  \item \textbf{The inverted pyramid.} An E2E-heavy suite is slow, flaky,
        and diagnostically mute --- a red test says ``something, somewhere''.
        Every E2E failure should be convertible into a missing lower-layer
        test; if it cannot, your layers have a hole.
  \item \textbf{Mocking what you do not own, naively.} A hand-built stub of
        a third-party API encodes your \emph{guess} about that API. The
        stub stays green while the vendor changes. Mitigations, in order:
        contract tests with the vendor's participation, record--replay
        cassettes re-recorded on a schedule, and at minimum a nightly
        smoke against the real sandbox.
  \item \textbf{Asserting on wall-clock.} \texttt{assert elapsed < 2.0}
        fails on a loaded CI runner and passes on an idle one --- the
        assertion measures the machine, not the code. Inject clocks and
        sleepers; assert on recorded virtual schedules. Reserve real
        latency assertions for scheduled runs on quiet hardware, and even
        then compare ratios (candidate vs.\ baseline), not absolutes.
  \item \textbf{Load-testing production without isolation.} A stress test
        against shared production is an outage you scheduled. Isolate:
        dedicated environment, synthetic tenants, synthetic data, and a
        blast-radius budget agreed with whoever gets paged.
  \item \textbf{Flaky retries masking real bugs.} ``Re-run and it passes''
        applied blindly converts race conditions into invisible inventory.
        Policy: two flakes in a week means fix or delete; automatic retry
        in CI is allowed only with failure-cause logging and a flake
        dashboard.
  \item \textbf{Fuzzing without shrinking or report hygiene.} A 4{,}000-key
        failing payload is not a bug report; shrink it. And always emit
        coverage of the campaign --- cases run, classes hit --- otherwise
        a crash-on-startup fuzzer reports ``0 failures'' forever.
  \item \textbf{Verifying contracts against a mock provider.} If the
        ``provider side'' of contract verification stubs the provider, the
        only thing verified is the pact file's internal consistency.
        Replay against the real application, in-process at minimum.
  \item \textbf{Provider states that seed the universe.} A state named
        ``standard test data'' that loads 400 fixtures hides which
        interaction needs what. Minimal, semantic, independent states ---
        or debugging a failure becomes archaeology.
  \item \textbf{Coverage floors as quality theater.} A 95\% line-coverage
        gate can be satisfied by tests that assert nothing. Floors guard
        against suite deletion; they do not manufacture confidence. Pair
        every floor with invariant-based tests that must fail when behavior
        breaks.
\end{enumerate}

\begin{warningbox}
The two most expensive words in API quality are ``it passed''. Passed
\emph{what}, running \emph{where}, against \emph{which} version of the
dependency, with \emph{which} seed? A gate whose inputs are not pinned
--- versions, seeds, pacts, specs --- manufactures confidence, not
evidence.
\end{warningbox}

%% ----------------------------------------------------------------------------
\section{Hands-On Production Lab \& Real-World Case Study}
\label{sec:ch19-lab}

\subsection{Lab: Run the Gate, Then Break the API}
\label{sec:ch19-labsteps}

Work from the book repository root in PowerShell. Estimated time: 60--90
minutes.

\begin{minted}{powershell}
cd code\ch19
pip install -r requirements.txt        # once
pytest -q                              # 15 tests: pact, pbt, fuzz, faults
\end{minted}

\begin{enumerate}
  \item \textbf{Baseline.} Run \texttt{pytest -q} (all green), then
        \texttt{python fuzz\_openapi\_mini.py} and note: 38 cases, zero
        failures on the healthy app; the buggy variant yields one failure
        shrunk to \texttt{price\_cents = 666}. Run
        \texttt{python load\_harness.py} and read the sweep: locate the
        knee level and compare the Little's Law estimate to the configured
        worker count.
  \item \textbf{Run the gate.} \texttt{python ci\_gate.py}; confirm all
        three sub-gates pass and \texttt{\$LASTEXITCODE} is 0. Read
        \texttt{pacts\textbackslash{}fulfillment-service-robot-registry.json}
        --- this is the contract the next step will break.
  \item \textbf{Break the contract.} In \texttt{ci\_gate.py}, change the
        healthy verifier factory from \texttt{create\_app} to
        \texttt{lambda: create\_app(breaking=True)} and re-run. The
        contract gate fails and names the missing field:
        \texttt{\$.sku: missing}. Restore the file.
  \item \textbf{Break the error contract.} Change the fuzz gate's healthy
        client to \texttt{create\_app(bug\_at\_price=666)} and re-run: the
        fuzz smoke fails with a 5xx report. Restore.
  \item \textbf{Break capacity.} In \texttt{ci\_gate.py}, drop the perf
        smoke to \texttt{rps=50, duration\_s=0.4} against
        \texttt{create\_workload\_app(work\_ms=400)} and observe queueing
        in a manual sweep; reason with Little's Law about which pool
        exhausts first.
  \item \textbf{Reflect.} For each break: which gate caught it, at which
        pyramid layer, and what would the same break have cost as a
        production incident?
\end{enumerate}

\subsection{War Story: The Deploy That Passed Everything Except the Contract}
\label{sec:ch19-warstory}

A fulfillment platform team (names changed; the failure is painfully
common) ran a genuinely good pipeline: 90\% coverage, integration
environments, canary deploys. The robot-registry team shipped what their
OpenAPI review called a ``cosmetic normalization'': the response field
\texttt{sku} became \texttt{part\_number}, matching the new ERP system's
vocabulary. Every provider-side test passed --- they had all been updated
in the same commit, which is exactly the problem. The staging integration
suite passed too, because the only consumer exercised there was the
provider team's own demo client.

The fulfillment service --- which read \texttt{sku} on every order line ---
did not fail at deploy time. It failed eleven minutes later, at the peak of
the hourly batch, as a tide of \texttt{KeyError: 'sku'} that the canary
analysis could not see because the canary received no fulfillment traffic.
Rollback took forty minutes; order backlog recovery took the weekend.

The postmortem's uncomfortable finding was that the consumer team
\emph{had} written a pact --- months earlier, in a spike --- and it sat in
a gist nobody owned. It was never published to the broker, the provider
team never knew it existed, and no \texttt{can-i-deploy} gate consumed it.
The failure was not technical; the loop of
Figure~\ref{fig:ch19-contract-loop} was never closed organizationally. The
fix was equally organizational: pact publication became a required artifact
of consumer CI, provider verification became a required step of provider
CI, and the deploy gate learned to answer ``which consumers have verified
this exact provider version?'' The technology in this chapter is the easy
half; wiring the artifact across the team boundary is the half that fails.

\begin{examplebox}
Mini-case, fully reproducible on your machine: the lab's step~3 is this
war story in miniature. The provider's own tests (Listing 19.1's app is
self-consistent in both modes) stay green under the rename; only the
consumer-recorded pact detects it. If you want the full emotional
experience, delete the pact file first and watch the gate turn green ---
then decide who on your team would have caught that in review.
\end{examplebox}

%% ----------------------------------------------------------------------------
\section{Self-Assessment Exercises \& Deep-Dive Questions}
\label{sec:ch19-exercises}

\begin{enumerate}[label=\textbf{E-\arabic*.}, leftmargin=1.9cm]
  \item \textbf{Layer audit.} Take the fifteen tests in
        \texttt{code\textbackslash{}ch19} and classify each into exactly
        one pyramid layer of Definition~\ref{def:test-layers}. Which layer
        is over-weighted relative to its defect-catching power, and why is
        that acceptable here?
  \item \textbf{Designated catcher.} For each defect --- (a) SKU validation
        regex wrong, (b) response field renamed, (c) database connection
        string mispointed, (d) retry backoff missing jitter, (e) p99
        regression after a serializer change --- name the cheapest pyramid
        layer that can catch it and sketch the test.
  \item \textbf{Contract authoring.} Add a third consumer,
        \texttt{billing-service}, that reads only \texttt{id} and
        \texttt{price\_cents}. Record its pact. Prove that renaming
        \texttt{name} breaks no consumer while renaming
        \texttt{price\_cents} breaks exactly one --- the usage-sensitivity
        schema diffing cannot express.
  \item \textbf{Provider-state discipline.} Refactor the seeded state into
        two independent states (``an empty registry'', ``a robot with id 7
        exists'') and replay. What breaks if two interactions share one
        app instance instead of one fresh app per interaction?
  \item \textbf{Metamorphic properties.} Extend Listing 19.4 with a
        metamorphic relation: creating the same payload twice (minus
        id) must yield equal records up to the \texttt{id} field. Why is
        this stronger than any single example assertion?
  \item \textbf{Coordinated omission.} Modify \texttt{open\_loop} to
        schedule by completion time (send next 1/rps after the previous
        finishes). Re-run the sweep at 400 rps and explain why the table
        now lies about the knee.
  \item \textbf{Blast radius.} Write the steady-state hypothesis and
        abort-criteria paragraph for running Listing 19.7's latency fault
        against a shared staging environment. What telemetry proves the
        experiment stayed in bounds?
  \item \textbf{Gate economics.} Your CI budget is four minutes. Allocate
        it across the gates of Listing 19.8 plus coverage, and defend the
        cuts. What moves to a nightly job, and what evidence do you lose?
\end{enumerate}

%% ----------------------------------------------------------------------------
\section{Interview Q\&A Vault}
\label{sec:ch19-interviews}

\begin{enumerate}[label=\textbf{Q\arabic*.}, leftmargin=1.7cm]
  \item \emph{[junior]} \textbf{What are the layers of the API test
        pyramid, and what does each catch?} Unit: pure logic (validators,
        serializers) --- logic bugs, in microseconds. Component: whole
        service in-process --- wiring, validation, error-shape bugs.
        Contract: consumer--provider agreement --- cross-team breakage.
        Integration: real infrastructure boundaries --- configuration and
        migration drift. E2E: golden journeys --- emergent whole-system
        failures only.
  \item \emph{[junior]} \textbf{Unit versus component test for an HTTP
        endpoint?} A unit test calls the validation function directly; a
        component test sends a real HTTP request through the real router
        and serializer, in-process. The component test catches URL
        mapping, content-negotiation, and error-shaping bugs the unit
        test cannot see.
  \item \emph{[junior]} \textbf{What is a consumer-driven contract?} An
        artifact authored by the \emph{consumer} recording the exact
        interactions it depends on, replayed by the \emph{provider's}
        build against real code. The direction of authorship is the whole
        point: it inverts who defines ``correct''.
  \item \emph{[junior]} \textbf{What is inside a pact file?} Consumer and
        provider names, and a list of interactions: request (method, path,
        body), expected response (status, the body subset actually read),
        and the named provider state required for replay.
  \item \emph{[junior]} \textbf{Soak versus stress versus spike?} Stress:
        ramp until failure to find the knee. Soak: expected load for
        hours to catch leaks and slow degradation. Spike: sudden 10x
        burst to test elasticity and cold caches.
  \item \emph{[junior]} \textbf{Why p99 and not average latency?} Averages
        hide tails; users experience tails. At 1\% bad requests and a
        million requests a day, ten thousand users had a bad day while the
        average looked fine.
  \item \emph{[mid]} \textbf{What makes a good provider state?} Minimal
        (arranges exactly what the interaction needs), independent (any
        order or subset), and semantic (named by domain meaning:
        ``a robot with id 7 exists'', not ``fixture set B''). States that
        seed the universe hide coupling and make failures undebuggable.
  \item \emph{[mid]} \textbf{Consumer-driven contracts versus OpenAPI
        schema diffing --- when each?} Contracts are usage-based: they
        catch breakage in fields consumers actually read, including
        semantic drift inside unchanged shapes, but cover only enrolled
        consumers. Schema diff is structural: it guards public APIs with
        unknown consumers but flags only syntactic incompatibility.
        Internal microservices with few consumers: Pact. Public API:
        diff plus both if possible.
  \item \emph{[mid]} \textbf{Fake versus stub versus mock?} Fake: working
        implementation with a shortcut (in-memory store, ASGI transport).
        Stub: canned answers. Mock: pre-programmed expectations that fail
        on unexpected calls. Prefer fakes for what you own; stubs at
        third-party edges; mocks sparingly, since they couple tests to
        interaction detail.
  \item \emph{[mid]} \textbf{When is mocking what you don't own
        acceptable?} At a hard third-party boundary, with canned responses
        derived from the vendor's documented contract --- and always paired
        with either a vendor-verified contract, a re-recorded cassette, or
        a scheduled sandbox smoke. The unforgivable version asserts on the
        mock's call list and calls it integration testing.
  \item \emph{[mid]} \textbf{What is property-based testing?} Instead of
        hand-picked examples, you define a generator (strategy) over the
        input space and an invariant that must hold for every generated
        input; the framework explores, and on failure shrinks to a minimal
        counterexample.
  \item \emph{[mid]} \textbf{Name three invariants every API should
        satisfy under generated inputs.} Never 5xx; every 2xx validates
        against the published response schema; every 4xx is RFC~9457
        problem-shaped.
  \item \emph{[mid]} \textbf{What is shrinking and why does it matter?}
        Automatic minimization of a failing input while the failure still
        reproduces. It converts ``500 on this 4{,}000-key blob'' into
        ``500 when \texttt{price\_cents} is 666'' --- the difference
        between noise and a bug report.
  \item \emph{[mid]} \textbf{How do you test timeouts and retries without
        sleeping?} Inject the sleeper. Production wires the real
        \texttt{sleep}; tests wire a virtual clock that records requested
        delays and advances instantly. Assert on the recorded backoff
        schedule (\texttt{[0.1, 0.2, 0.4]}), which is the real behavioral
        contract, in microseconds of wall time.
  \item \emph{[mid]} \textbf{How do you test a circuit breaker
        deterministically?} Scripted fault injection indexed by request
        number (not randomness): $N$ consecutive exhausted calls, then
        assert the $(N{+}1)$-th call fails fast by counting requests that
        reached the dependency. Then inject one success and assert the
        breaker resets.
  \item \emph{[mid]} \textbf{Risks of record--replay cassettes?} Secrets
        and PII baked into fixtures (sanitize); silent rot as the real
        service evolves (re-record on schedule, or at least smoke against
        reality); and false confidence because replay cannot diverge ---
        it can only be stale.
  \item \emph{[mid]} \textbf{What belongs in a CI quality gate for an
        API?} Fast hermetic layers: unit + property suites with a coverage
        floor, contract verification per published pact, OpenAPI diff
        versus the released spec, a bounded fuzz smoke, and a structural
        perf smoke. Deep load and soak move to scheduled jobs.
  \item \emph{[mid]} \textbf{What is the value --- and the failure mode ---
        of coverage floors?} Value: they make suite deletion visible and
        force new code to arrive with \emph{some} test. Failure mode: they
        reward assertion-free tests that touch lines. Pair floors with
        invariant-based tests and review test \emph{assertions}, not just
        diffs.
  \item \emph{[mid]} \textbf{How do you handle flaky tests?} Track flake
        rate per test; two flakes in a week means fix or delete;
        quarantine with an SLA, never as a graveyard; make CI retries log
        failure causes; prefer determinism (derandomized seeds, virtual
        clocks) over retry budgets.
  \item \emph{[senior]} \textbf{Open versus closed workload models, and
        what is coordinated omission?} Closed: $N$ workers each wait for
        the response before sending again --- offered load shrinks when the
        system slows. Open: fixed arrival rate regardless of completions.
        Coordinated omission is the closed-model pathology where load
        scheduled \emph{during} a latency spike is silently skipped,
        producing healthy-looking results for an unhealthy system. Latency
        and saturation work demands open models.
  \item \emph{[senior]} \textbf{Derive Little's Law and apply it to
        sizing.} Over a long window, the area under $L(t)$ equals the sum
        of residence times; dividing by $T$ gives $L = \lambda W$.
        Sizing: at 800 req/s and 30 ms mean latency, $L = 24$ requests are
        resident --- every pool and concurrency limit must exceed 24 with
        burst headroom, or queueing begins.
  \item \emph{[senior]} \textbf{What is the knee point and how do you find
        it operationally?} The offered load where latency stops growing
        linearly and turns hyperbolic ($W \approx S/(1-\rho)$). Find it
        with an open-model sweep: fixed windows at increasing RPS, plotting
        p99 and error rate; the knee is the first level where p99 exceeds
        roughly 10x baseline or errors appear. Provision at 60--70\% of
        the knee.
  \item \emph{[senior]} \textbf{How does Schemathesis-style fuzzing differ
        from naive random fuzzing?} It is schema-guided: candidates are
        drawn from the OpenAPI constraints (\texttt{minLength}$-1$,
        \texttt{maxLength}$+1$, \texttt{minimum}$-1$, type confusions,
        missing/extra keys), so it concentrates on validation boundaries
        where bugs live, and it checks contract invariants (never 5xx,
        schema-valid 2xx, problem-shaped 4xx) instead of merely watching
        for crashes.
  \item \emph{[senior]} \textbf{What does stateful rule-based testing
        catch that single-call properties cannot?} Sequence-dependent
        bugs: lost updates, read-your-writes violations, idempotency-key
        reuse corruption, cross-resource leakage. You generate random
        operation interleavings, maintain a local model, and assert
        agreement after every step.
  \item \emph{[senior]} \textbf{Define steady-state hypothesis and blast
        radius.} Steady state: a measurable behavioral definition of
        ``fine'' (p99 and error rate over a window), never an internal
        proxy. Hypothesis: the system stays in steady state under a
        specified fault. Blast radius: the maximal scope the experiment can
        harm --- start at one in-process middleware, widen only as
        confidence and abort automation mature.
  \item \emph{[senior]} \textbf{Give the fault taxonomy for HTTP APIs and
        one resilience defect each class reveals.} Latency (missing
        timeouts), synthesized 5xx (retry storms without jitter), abort
        (unhandled exception paths), partition (failed failover), resource
        exhaustion (queue growth, priority inversion).
  \item \emph{[senior]} \textbf{Performance regression budgets in CI when
        runners are noisy?} Do not assert wall-clock latency on shared
        runners: assert structure (completion, zero 5xx) in the PR gate,
        and run ratio-based comparisons (candidate versus same-run
        baseline) on dedicated quiet hardware as a scheduled job with
        archived histograms.
  \item \emph{[senior]} \textbf{A provider's tests all pass, the deploy
        ships, consumers break. What process failed?} The provider tested
        against itself. No consumer contract was verified for the deployed
        version: pact publication, provider-side replay, and the
        \texttt{can-i-deploy} gate were missing or not blocking. Schema
        diff may also have missed it if the change was semantic.
  \item \emph{[staff]} \textbf{Design a load test that finds the knee of a
        new checkout API.} Model the workload (open model; realistic
        request mix and payload sizes); isolate an environment; baseline
        at trivial load; sweep offered RPS in fixed short windows past
        predicted capacity ($c/S$); record p50/p99/error rate per level;
        flag the knee at 10x-baseline p99 or first error; cross-check with
        Little's Law ($L = \lambda W$ versus pool sizes); set the budget
        at 60--70\% of the knee; convert the sweep into a scheduled
        regression job and a structural PR smoke.
  \item \emph{[staff]} \textbf{How would you roll out consumer-driven
        contracts across forty microservices?} Start with the three
        highest-change-frequency provider/consumer pairs; stand up a
        broker; make publication a CI artifact and verification a provider
        CI step; run gates advisory-first for a month, publish the
        catch statistics, then make them blocking; invest in
        provider-state libraries per domain; measure pact coverage of
        production traffic, not of repositories.
  \item \emph{[staff]} \textbf{What does the delivery-performance evidence
        (Accelerate, SRE practice) say about test automation?} Fast,
        trusted, automated gates correlate with \emph{both} deployment
        throughput and stability --- quality gates are a delivery feature.
        The organizational failure mode is slow, flaky suites that train
        engineers to bypass them; the fix is engineering discipline in the
        gates themselves, not more process around them.
  \item \emph{[staff]} \textbf{Your pbt suite passes for months, then
        fails on a seed nobody can reproduce. Diagnose.} Check for hidden
        nondeterminism: unseeded randomness, wall-clock reads, dict/set
        ordering, environment dependence, leaked state between examples,
        or a live dependency sneaking into the ``offline'' suite. Then
        pin everything (derandomize, inject clocks, fresh app per test)
        and add the failing example as a regression test.
\end{enumerate}

%% ----------------------------------------------------------------------------
\section{External Bibliography \& Further Reading}
\label{sec:ch19-bibliography}

\begin{itemize}
  \item \textbf{Pact documentation} --- \emph{Pact: Consumer-Driven
        Contracts}, \texttt{https://docs.pact.io}~\cite{pactdocs}. The
        reference for the consumer/pact-file/broker/verification loop,
        provider states, and \texttt{can-i-deploy}.
  \item \textbf{OpenAPI 3.1 Specification} --- the schema source our
        fuzzer consumes; pair with any OpenAPI-diff tooling for structural
        compatibility gates~\cite{openapi31}.
  \item \textbf{Steve Freeman \& Nat Pryce}, \emph{Growing
        Object-Oriented Software, Guided by Tests} (Addison-Wesley, 2009).
        The canonical treatment of mock-based design, the doubles
        taxonomy in practice, and why tests should express behavior.
  \item \textbf{Gerard Meszaros}, \emph{xUnit Test Patterns: Refactoring
        Test Code} (Addison-Wesley, 2007). The precise dummy/stub/fake/
        spy/mock vocabulary used in Section~\ref{sec:ch19-doubles}.
  \item \textbf{Hypothesis documentation} ---
        \texttt{https://hypothesis.readthedocs.io}. Strategies, shrinking,
        \texttt{derandomize}, and \texttt{RuleBasedStateMachine} for
        stateful API testing.
  \item \textbf{Schemathesis documentation} ---
        \texttt{https://schemathesis.readthedocs.io}. Production-grade
        OpenAPI-driven fuzzing whose miniature we built in Listing 19.5.
  \item \textbf{k6 documentation} --- \texttt{https://grafana.com/docs/k6}.
        Arrival-rate (open-model) executors, thresholds, and checks;
        Locust (\texttt{https://docs.locust.io}) is the Python-native
        alternative.
  \item \textbf{Betsy Beyer et al.} (eds.), \emph{Site Reliability
        Engineering} (O'Reilly, 2016)~\cite{beyer2016sre} --- the testing
        and release-engineering chapters, plus SLO thinking for load-test
        budgets.
  \item \textbf{Michael T. Nygard}, \emph{Release It!, 2nd ed.} (Pragmatic
        Bookshelf, 2018)~\cite{nygard2018releaseit} --- stability patterns
        and anti-patterns; the conceptual home of circuit breakers and
        capacity disasters.
  \item \textbf{Martin Kleppmann}, \emph{Designing Data-Intensive
        Applications} (O'Reilly, 2017)~\cite{kleppmann2017ddia} ---
        partial failure, partitions, and why distributed systems need
        fault injection.
  \item \textbf{Nicole Forsgren, Jez Humble, Gene Kim},
        \emph{Accelerate} (IT Revolution, 2018) --- the empirical link
        between fast, trusted test automation and delivery performance.
  \item \textbf{FastAPI and httpx documentation}~\cite{fastapidocs,httpxdocs}
        --- application factories, ASGI transports, and
        \texttt{MockTransport} interception patterns.
  \item \textbf{RFC 9457}, \emph{Problem Details for HTTP APIs} --- the
        error shape our invariants assert on every 4xx.
  \item \textbf{Gil Tene}, ``How NOT to Measure Latency'' (talk, 2011) ---
        the definitive coordinated-omission rant behind our open-model
        insistence.
\end{itemize}

%% ----------------------------------------------------------------------------
\section{Capstone Projects}
\label{sec:ch19-capstones}

\subsection*{Capstone 19.1 [junior] --- A Second Consumer}
\textbf{Objective:} Extend the chapter's contract suite with a new consumer
to see usage-based verification in action.

\textbf{Inputs/Datasets:} \texttt{code\textbackslash{}ch19\textbackslash{}}
(robot registry, \texttt{pact\_mini.py}); synthetic Northwind Robotics data
only.

\textbf{Steps:}
\begin{enumerate}
  \item Invent \texttt{billing-service}, which reads only \texttt{id} and
        \texttt{price\_cents} from \texttt{GET /robots/\{id\}}.
  \item Record its pact with its own provider state
        (``a robot with id 11 exists'').
  \item Verify both pacts against the healthy app, then against
        \texttt{create\_app(breaking=True)}.
  \item Introduce a third variant that renames \texttt{name} only, and show
        no consumer pact fails.
\end{enumerate}

\textbf{Acceptance Criteria:}
\begin{itemize}
  \item[$\square$] Two pact files exist and verify green against the
        healthy app.
  \item[$\square$] The \texttt{sku} rename fails exactly one pact; the
        \texttt{name} rename fails none.
  \item[$\square$] Reports name the failing field with a JSONPath-style
        location.
\end{itemize}

\textbf{Stretch Goals:} Add a \texttt{can-i-deploy(provider\_version)}
function consulting a local JSON ``broker'' mapping versions to verified
pacts.

\subsection*{Capstone 19.2 [mid] --- Property Suite for a New Resource}
\textbf{Objective:} Bring a second resource under property-based and fuzz
coverage.

\textbf{Inputs/Datasets:} \texttt{code\textbackslash{}ch19\textbackslash{}};
a new \texttt{/warehouses} resource you add to the registry (fields:
\texttt{code}, \texttt{city}, \texttt{capacity\_m2}).

\textbf{Steps:}
\begin{enumerate}
  \item Implement the resource with RFC~9457 error handling.
  \item Write hostile-payload Hypothesis tests encoding the three universal
        invariants.
  \item Extend the fuzzer's schema excerpt and candidate generation for the
        new fields.
  \item Plant two distinct bugs (a 500 trigger; a schema-violating 201) and
        confirm both are found and shrunk.
\end{enumerate}

\textbf{Acceptance Criteria:}
\begin{itemize}
  \item[$\square$] \texttt{derandomize=True} suites pass repeatably.
  \item[$\square$] Both planted bugs produce minimal shrunk reproducers.
  \item[$\square$] Zero 5xx and zero non-problem 4xx across the campaign.
\end{itemize}

\textbf{Stretch Goals:} Add a rule-based stateful machine over
create/read/update sequences for warehouses, with a local model oracle.

\subsection*{Capstone 19.3 [mid] --- Record--Replay Cassette Library}
\textbf{Objective:} Build a minimal VCR for \texttt{httpx} and learn where
cassettes rot.

\textbf{Inputs/Datasets:} \texttt{code\textbackslash{}ch19\textbackslash{}};
a JSON cassette format you design (request key $\rightarrow$ recorded
status/headers/body).

\textbf{Steps:}
\begin{enumerate}
  \item Implement a recording transport that writes cassettes and a
        replaying transport that serves them.
  \item Add sanitization hooks (strip \texttt{Authorization}, redact
        fields).
  \item Record the robot-registry interactions; replay with the app
        offline.
  \item Mutate the app (change a field) and demonstrate that replay cannot
        detect it --- then add a cassette ``freshness date'' check.
\end{enumerate}

\textbf{Acceptance Criteria:}
\begin{itemize}
  \item[$\square$] Replay works with zero application code.
  \item[$\square$] No credential-like header ever reaches disk.
  \item[$\square$] A stale cassette is detected and reported by date or
        hash.
\end{itemize}

\textbf{Stretch Goals:} Request matching on method+path+canonical body;
strict versus lax matching modes.

\subsection*{Capstone 19.4 [senior] --- Capacity Report from a Knee Sweep}
\textbf{Objective:} Produce a capacity analysis for a service using the
harness, defended with Little's Law.

\textbf{Inputs/Datasets:} \texttt{code\textbackslash{}ch19\textbackslash{}}
load harness; a workload app with tunable \texttt{work\_ms} and
\texttt{max\_concurrent}.

\textbf{Steps:}
\begin{enumerate}
  \item Predict capacity analytically: $c/S$, with $\rho \to 1$ latency
        shape.
  \item Run open-model sweeps at three \texttt{work\_ms} settings; tabulate
        knee points; compare with prediction.
  \item Cross-check with a closed-model run and $L = \lambda W$.
  \item Write the one-page report: knee, recommended provisioning headroom,
        and the perf budget to encode in CI.
\end{enumerate}

\textbf{Acceptance Criteria:}
\begin{itemize}
  \item[$\square$] Predicted and measured knees agree within a factor of
        two, with the discrepancy explained.
  \item[$\square$] The report includes a knee table, a Little's Law
        cross-check, and an explicit headroom recommendation.
  \item[$\square$] No wall-clock assertion anywhere; numbers labeled
        illustrative.
\end{itemize}

\textbf{Stretch Goals:} Add a coordinated-omission buggy scheduler variant
and quantify how much it overstates capacity.

\subsection*{Capstone 19.5 [staff] --- Quality-Gate Architecture for a Fleet}
\textbf{Objective:} Design the gate architecture for a (fictional)
ten-service Northwind Robotics fleet.

\textbf{Inputs/Datasets:} The chapter's gate; a written fleet topology you
invent (services, consumers, deploy cadences).

\textbf{Steps:}
\begin{enumerate}
  \item Allocate pyramid layers per service with designated catchers per
        defect class.
  \item Specify the broker topology: pact storage, version tagging,
        \texttt{can-i-deploy} semantics per environment.
  \item Budget CI minutes per gate; move deep sweeps and soaks to
        scheduled jobs with archived evidence.
  \item Write the flake policy (quarantine SLA, deletion rule, dashboards)
        and the advisory-to-blocking rollout plan for contract gates.
\end{enumerate}

\textbf{Acceptance Criteria:}
\begin{itemize}
  \item[$\square$] Every defect class in your taxonomy has exactly one
        designated cheapest catcher.
  \item[$\square$] The gate budget fits four minutes per PR with a written
        justification for each cut.
  \item[$\square$] The rollout plan includes measurable adoption metrics
        (pact coverage of production traffic) and a self-check requirement
        for every gate.
\end{itemize}

\textbf{Stretch Goals:} Add mutation testing of the gates themselves
(deliberately broken variants) as a scheduled meta-gate.

%% ----------------------------------------------------------------------------
\section{Python Dependencies Appendix}
\label{sec:ch19-deps}

\begin{minted}{text}
# code\ch19\requirements.txt -- Chapter 19, Python 3.11+, offline-first
pytest>=8.0
hypothesis>=6.100
fastapi>=0.110
httpx>=0.27
pydantic>=2.7
jsonschema>=4.21
\end{minted}

Installation (PowerShell, from the track folder):

\begin{minted}{powershell}
python -m venv .venv
.\.venv\Scripts\Activate.ps1
pip install -r code\ch19\requirements.txt
pytest code\ch19 -q
\end{minted}

\subsection{pytest}
\textbf{Description:} The test runner and fixture framework used by every
suite in the chapter.
\textbf{Why included:} Discovery, \texttt{tmp\_path} fixtures, rich
assertion introspection, and exit codes compose into CI gates without
glue code.
\textbf{Alternatives:} \texttt{unittest} (stdlib, heavier boilerplate);
\texttt{nose2} (largely dormant).
\textbf{Installation:} \texttt{pip install "pytest>=8.0"}.

\subsection{hypothesis}
\textbf{Description:} Property-based testing library: strategies,
shrinking, and rule-based state machines.
\textbf{Why included:} Powers Listing 19.4's invariant suites;
\texttt{derandomize=True} gives deterministic CI replays.
\textbf{Alternatives:} \texttt{schemathesis} (OpenAPI-driven API fuzzing
built on Hypothesis --- use it when you may add a dependency per spec);
hand-rolled generators (Listing 19.5 shows the floor of that path).
\textbf{Installation:} \texttt{pip install "hypothesis>=6.100"}.

\subsection{fastapi}
\textbf{Description:} The ASGI framework implementing the robot registry
under test~\cite{fastapidocs}.
\textbf{Why included:} Real routing, validation, and serialization give
component and contract tests something true to verify.
\textbf{Alternatives:} Starlette (lower level, same ASGI base); Flask
(WSGI --- needs a different in-process testing story).
\textbf{Installation:} \texttt{pip install "fastapi>=0.110"}.

\subsection{httpx}
\textbf{Description:} HTTP client with both sync and async APIs and,
crucially, transport injection~\cite{httpxdocs}.
\textbf{Why included:} \texttt{ASGITransport} runs everything in-process;
\texttt{MockTransport} underpins stub-based client tests; custom
\texttt{AsyncBaseTransport} hosts the resilient transport of Listing 19.7.
\textbf{Alternatives:} \texttt{respx} (declarative route mocking on top of
httpx); \texttt{requests} + \texttt{responses} (sync-only, older idiom).
\textbf{Installation:} \texttt{pip install "httpx>=0.27"}.

\subsection{pydantic}
\textbf{Description:} Validation and serialization models behind FastAPI.
\textbf{Why included:} Field constraints generate the 422 boundary our
fuzzer explores; \texttt{model\_json\_schema()} is the bridge between code
and contract.
\textbf{Alternatives:} \texttt{dataclasses} + manual validation (stdlib,
no schema generation); \texttt{attrs} + \texttt{cattrs}.
\textbf{Installation:} \texttt{pip install "pydantic>=2.7"}.

\subsection{jsonschema}
\textbf{Description:} JSON Schema validation for Python.
\textbf{Why included:} Asserts the ``2xx validates against the published
schema'' invariant independently of the server's own serializer --- the
checker must not share the checker's bugs.
\textbf{Alternatives:} \texttt{pydantic.TypeAdapter} (faster, but shares
implementation lineage with the app --- weaker independence);
\texttt{fastjsonschema}.
\textbf{Installation:} \texttt{pip install "jsonschema>=4.21"}.

\paragraph{Notable tools intentionally not required.}
\texttt{schemathesis} (we build its core idea by hand for instruction),
\texttt{pact-python} (we implement the loop explicitly), \texttt{locust}
and \texttt{k6} (we implement open/closed models in 100 lines),
\texttt{respx} (MockTransport suffices), and \texttt{pytest-asyncio}
(we call \texttt{asyncio.run} directly to keep the dependency surface at
six packages).

%% ----------------------------------------------------------------------------
\section{Chapter Summary \& Key Takeaways}
\label{sec:ch19-summary}

\begin{itemize}
  \item The test pyramid is portfolio allocation: every defect class gets a
        \emph{designated cheapest catcher}; E2E stays thin because it is
        slow, flaky, and diagnostically mute.
  \item Consumer-driven contracts invert correctness: consumers record what
        they rely on, providers replay it against real code with minimal
        named states, and \texttt{can-i-deploy} turns verification into a
        release gate~\cite{pactdocs}. Schema diffing guards public
        surfaces; contracts guard known consumers; use both.
  \item Property-based testing and schema-driven fuzzing buy coverage of
        input spaces you cannot enumerate; the universal invariants are
        never-5xx, schema-valid 2xx, RFC~9457 4xx; shrinking turns crashes
        into bug reports; stateful rules catch sequence bugs.
  \item Test doubles have precise meanings; prefer fakes for what you own,
        stubs at third-party edges, sanitized cassettes with freshness
        dates; never let a mock of someone else's API be your integration
        evidence~\cite{httpxdocs}.
  \item Determinism is engineered: injected sleepers and virtual clocks
        test retries, timeouts, and breakers in microseconds; scripted
        fault injection makes chaos reproducible in CI.
  \item Load testing is the search for the knee: open models avoid
        coordinated omission; Little's Law $L = \lambda W$ connects
        throughput, latency, and concurrency; provision at 60--70\% of the
        knee; load/stress/soak/spike answer different questions
        \cite{nygard2018releaseit,beyer2016sre}.
  \item Quality gates are code: hermetic, fast, deterministic, legible,
        and \emph{self-checking} --- every gate in this chapter proves it
        can fail by running against a deliberately broken variant.
\end{itemize}

%% ----------------------------------------------------------------------------
\section{Quick-Reference Card}
\label{sec:ch19-quickref}

\begin{table}[htbp]
\centering\small
\begin{tabularx}{\textwidth}{l X X}
\toprule
\textbf{Layer} & \textbf{Catches} & \textbf{Tooling (this chapter)} \\
\midrule
Unit        & validation/serialization logic & pytest on pure functions \\
Component   & routing, wiring, error shapes & httpx + ASGI transport (in-process) \\
Contract    & cross-team breakage & pact recorder + verifier, OpenAPI diff \\
Integration & config, auth, migration drift & real infra in ephemeral env \\
E2E         & emergent whole-system failures & a handful of golden journeys \\
Property/fuzz & input-space and sequence bugs & Hypothesis, mini OpenAPI fuzzer \\
Load        & capacity, knee, regression & asyncio harness, budgets \\
Chaos       & resilience claims & scripted fault middleware \\
\bottomrule
\end{tabularx}
\caption{Test-layer responsibility matrix.}
\label{tab:ch19-layers}
\end{table}

\begin{table}[htbp]
\centering\small
\begin{tabularx}{\textwidth}{l X}
\toprule
\textbf{Fault} & \textbf{Assertion pattern in tests} \\
\midrule
Latency   & virtual clock recorded the delay; no wall-clock assertion \\
Error 5xx & retried with recorded backoff schedule; breaker trips at threshold \\
Abort     & surfaced as transport error after bounded retries \\
Partition & failover/fallback engaged; steady state held \\
Exhaustion & queue bounds respected; load-shedding, not collapse \\
\bottomrule
\end{tabularx}
\caption{Fault taxonomy and how to assert on it deterministically.}
\label{tab:ch19-faults}
\end{table}

\begin{table}[htbp]
\centering\small
\begin{tabularx}{\textwidth}{l X}
\toprule
\textbf{Quantity} & \textbf{Rule} \\
\midrule
Little's Law & $L = \lambda W$; size every pool above $\lambda W$ with burst headroom \\
Utilization & $\rho = \lambda S / c$; latency $\sim S/(1-\rho)$ as $\rho \to 1$ \\
Knee & first sweep level with p99 $>$ 10$\times$ baseline p50, or any errors \\
Provisioning & run production at 60--70\% of the knee \\
Perf in CI & structural assertions on PRs; ratio budgets on quiet scheduled runs \\
\bottomrule
\end{tabularx}
\caption{Little's Law and saturation cheat sheet.}
\label{tab:ch19-little}
\end{table}

\section{Standardized Evidence Addendum}
\subsection{Benchmark Snapshot Template}
\begin{table}[h]
  \centering
  \small
  \begin{tabularx}{\textwidth}{l c c c c}
    \toprule
    \textbf{Config} & \textbf{p50 latency} & \textbf{p99 latency} & \textbf{Error rate} & \textbf{Throughput} \\
    \midrule
    Baseline & -- & -- & -- & 1.00x \\
    Candidate A & -- & -- & -- & -- \\
    Candidate B & -- & -- & -- & -- \\
    \bottomrule
  \end{tabularx}
  \caption{Minimum benchmark table to keep per chapter.}
\end{table}

\subsection{Request Trace Template}
\begin{itemize}
  \item \textbf{Request:} one representative API call for this chapter's topic.
  \item \textbf{Before:} behavior (headers, payload, latency, failure mode) with the naive approach.
  \item \textbf{After:} behavior with the technique in this chapter.
  \item \textbf{Result:} what changed in correctness, resilience, latency, and operability.
\end{itemize}

\subsection{Cost and Latency Budget Snapshot}
\begin{table}[h]
  \centering
  \small
  \begin{tabularx}{\textwidth}{l c c}
    \toprule
    \textbf{Stage} & \textbf{Latency budget} & \textbf{Cost driver} \\
    \midrule
    TLS / connection setup & -- & handshake round trips, keep-alive hit rate \\
    Request validation & -- & schema complexity, CPU per request \\
    Business logic and I/O & -- & downstream fan-out, query cost \\
    Serialization and egress & -- & payload size, compression ratio \\
    \bottomrule
  \end{tabularx}
  \caption{Operational budget template for this chapter's technique.}
\end{table}

\section{Configuration Preset Template}
\begin{minted}{text}
# api_policy.yaml
policy_version: "v1.0.0"
component: "<chapter_topic>"
timeout_ms: 0
retry_budget: 0
fallback_strategy: "fail_fast"
rollback_trigger: "error_budget_burn_gt_threshold"
\end{minted}

\section{CI Quality Gate Specification}
\begin{itemize}
  \item contract tests green against the pinned OpenAPI/AsyncAPI/Protobuf spec,
  \item no breaking schema changes without a version bump,
  \item error responses conform to RFC 9457 problem-details shape,
  \item benchmark deltas within approved regression bounds (latency and throughput),
  \item policy version and rollback plan present in release metadata.
\end{itemize}

\section{Incident Runbook Template}
\begin{enumerate}
  \item Detect regression from SLO burn-rate alerts or consumer reports.
  \item Scope impacted endpoints, versions, and consumer segments.
  \item Contain impact (rollback, feature flag, rate-limit tightening, or traffic shift).
  \item Diagnose with traces, structured logs, and contract diffs.
  \item Recover with validated fix and verified rollout procedure.
  \item Prevent recurrence via new regression tests and CI gates.
\end{enumerate}

\section{Cross-Chapter Decision Card}
\begin{table}[h]
  \centering
  \small
  \begin{tabularx}{\textwidth}{l X X}
    \toprule
    \textbf{Dimension} & \textbf{Choose this approach when} & \textbf{Prefer an alternative when} \\
    \midrule
    Use case & -- & -- \\
    Latency budget & -- & -- \\
    Operational cost & -- & -- \\
    Team maturity & -- & -- \\
    \bottomrule
  \end{tabularx}
  \caption{Decision card to complete per chapter.}
\end{table}

\section{ROI Model and Build-vs-Buy}
\begin{itemize}
  \item \textbf{Build when:} the capability is differentiating, compliance forbids vendors, or scale makes vendor pricing irrational.
  \item \textbf{Buy when:} the capability is commodity (auth, gateways, observability), time-to-market dominates, or staffing is thin.
  \item \textbf{Hidden costs to model:} on-call load, upgrade cadence, expertise ramp, data egress fees.
\end{itemize}

\section{Disaster Recovery Notes (RPO/RTO)}
\begin{itemize}
  \item Define recovery point objective (RPO): maximum acceptable data loss window.
  \item Define recovery time objective (RTO): maximum acceptable restoration time.
  \item Test restores regularly; an untested backup is not a backup.
\end{itemize}

